%-----------------------------------------------------------------------------
% Załącznik V
%-----------------------------------------------------------------------------

\cleardoublepage
\section*{V. Piotrkowicze i~Okrzejowie z~Dębego}
\label{sec:piotrkowicze}

Prapoczątków rodzin Piotrkowiczów i~Okrzejów należy szukać w~XVII~wieku
w~Starostwie Dębe, w~ziemi czerskiej województwa mazowieckiego, które było
wówczas częścią ziem królewskich II~Rzeczypospolitej. Obie te~rodziny zostały
wymienione jako chłopi zamieszkujący wieś Dębe w~\enquote{Inwentarzu dóbr
dzierżawy Dembe} \cite{debe1}~sporządzonym w~1772~roku przez rotmistrza
Andrzeja Olędzkiego na~polecenie Komisji Skarbowej Koronnej (patrz: ryc.
\ref{fig:okrzejowie_1772} oraz \ref{fig:piotrkowicze_1772}).

\begin{figure}[!ht]
    \vspace*{0.5cm}
    \centering \includegraphics[width=0.9\linewidth]{
        Dembe_Okrzejowie.jpg}
    \captionsetup{format=hang}
    \caption{Fragment inwentarza dóbr dzierżawy Dębe z~1772~roku -- na
    czerwono podkreślono wpis dotyczący rodziny Andrzeja Okrzei \cite{debe1}.}
    \label{fig:okrzejowie_1772}
\end{figure}

Liczne wpisy dotyczące rodzin Okrzejów i~Piotrkowiczów można również odnaleźć
w~aktach metrykalnych parafii pod wezwaniem Narodzenia Najświętszej Maryi
Panny w~Mińsku Mazowieckim \cite{par_minsk0}, do~której należała wówczas wieś
Dębe. Stąd też wiadomo, iż~Andrzej
Okrzeja\indeksosoby{Andrzej}{Okrzeja}{1721}{Dębe Wielkie} (ur.~1721~r.~--
zm.~1778~r.), wymieniony w~inwentarzu jako ojciec rodziny Okrzejów, urodził
się 5~kwietnia 1721~roku w~Dębem jako syn Andrzeja
\indeksosoby{Andrzej}{Okrzeja}{1765}{Dębe Wielkie} (ur.~1665~r.~-- zm.~1739~r.)
i~Krystyny\indeksosoby[Garwolak]{Krystyna}{Okrzeja}{1686}{Dębe Wielkie} z~domu
Garwolak (ur.~1686~r.~-- zm.~1746~r.). W~1746~roku ożenił się
z~Katarzyną
Brzost\indeksosoby[Brzost]{Katarzyna}{Okrzeja}{1719}{Brzóze} (ur.~1719~r.~--
zm.~1793~r.), z~którą dochował się co~najmniej ośmiorga dzieci:
Wojciecha\indeksosoby{Wojciech}{Okrzeja}{1747}{Dębe Wielkie} (ur.~1747~r.~--
zm. 1820~r.), Franciszka\indeksosoby{Franciszka}{Okrzeja}{1748}{Dębe Wielkie}
(ur.~1748~r.~-- zm. 1770~r.), Józefa\indeksosoby{Józef}{Okrzeja}{1749}{Dębe Wielkie}
(ur.~1749~r.~-- zm.~?), Szczepana\indeksosoby{Szczepan}{Okrzeja}{1750}{Dębe Wielkie}
(ur.~1750~r.~-- zm.~?), Macieja\indeksosoby{Maciej}{Okrzeja}{1752}{Dębe Wielkie}
(ur.~1752~r.~-- zm. 1813~r.), Mateusza\indeksosoby{Mateusz}{Okrzeja}{1754}{Dębe Wielkie}
(ur.~1754~r.~-- zm.~?), Marianny\indeksosoby{Marianna}{Okrzeja}{1757}{Dębe Wielkie}
(ur.~1757~r.~-- zm.~?) oraz Urszuli\indeksosoby{Urszula}{Okrzeja}{1762}{Dębe Wielkie}
(ur.~1762~r.~-- zm.~?). Inwentarz z~1772~roku wspomina natomiast jedynie
o~piątce dzieci Andrzeja (pisownia oryginalna):

\begin{centeredquote}
\textit{Andrzey Okrzeia Synowie Woytek Zonaty Maciey Lat 18. Mateusz Lat 10.
Corki Marya Lat 11. Orszula Lat 7.}
\end{centeredquote}

\begin{figure}[!ht]
    \vspace*{-0.5cm}
    \centering \includegraphics[width=1.0\linewidth]{
        Dembe_Piotrkowicze.jpg}
    \captionsetup{format=hang}
    \caption{Fragment inwentarza dóbr dzierżawy Dębe z~1772~roku -- na
    czerwono podkreślono wpis dotyczący rodziny Bartłomieja Piotrkowicza
    \cite{debe1}.}
    \label{fig:piotrkowicze_1772}
\end{figure}

\begin{figure}[!ht]
    \centering \includegraphics[width=1.0\linewidth]{
        1759_Bartłomiej_Piotrkowicz_Zofia_Kotczyk_akt_ślubu_parafia_Mińsk_Mazowiecki_wpis_67.jpg}
    \captionsetup{format=hang}
    \caption{Akt ślubu Bartłomieja Piotrkowicza i~Zofii Kotczyk -- par.
    Mińsk Mazowiecki 1759~rok (67/1759) 
    \cite{par_minsk0}.}
    \label{fig:bpiotrkowicz_1759}
\end{figure}

Jeżeli chodzi natomiast o~rodzinę Piotrkowiczów to akta metrykalne parafii
w~Mińsku Mazowieckim wskazują, iż~Bartłomiej
Piotrkowicz\indeksosoby{Bartłomiej}{Piotrkowicz}{1741}{Dębe Wielkie}
(ur.~1741~r.~-- zm.~1791~r.) urodził się w~1741~roku w~Dębem, w~1759~roku ożenił
się~z~Zofią Kotczyk\indeksosoby[Kotczyk]{Zofia}{Piotrkowicz}{1743}{Dębe
Wielkie} (ur.~1743~r.~-- zm.~1793~r.) panną pochodzącą z~rodzinnej wsi
Piotrkowiczów (patrz: ryc. \ref{fig:bpiotrkowicz_1759}). Para dochowała się
wspólnie co~najmniej piątki dzieci:

\begin{itemize}
\item \textbf{Pawła}\indeksosoby{Paweł}{Piotrkowicz}{1760}{Dębe Wielkie}
(ur.~1760~r.~-- zm.~?) -- patrz: ryc. \ref{fig:ppiotrkowicz_1760},
\item \textbf{Kacpra}\indeksosoby{Kacper}{Piotrkowicz}{1762}{Dębe Wielkie}
(ur.~1762~r.~-- zm.~?) -- patrz: ryc. \ref{fig:kpriotrkowicz_1762},
\item Marianny\indeksosoby{Marianna}{Piotrkowicz}{1766}{Dębe Wielkie}
(ur.~1766~r.~-- zm.~?),
\item \textbf{Sebastiana}\indeksosoby{Sebastian}{Piotrkowicz}{1768}{Dębe Wielkie}
(ur.~1768~r.~-- zm.~1818~r.) -- patrz: ryc. \ref{fig:spiotrkowicz_1768},
\item Agnieszki\indeksosoby{Agnieszka}{Piotrkowicz}{1768}{Dębe Wielkie}
(ur.~1768~r.~-- zm.~?).
\end{itemize}

\enquote{Inwentarz dóbr dzierżawy Dembe} z~1772~roku
wspomina natomiast jedynie o~czwórce dzieci Bartłomieja (pisownia oryginalna):

\begin{centeredquote}
    \textit{Bartek Piotrkowicz Paweł Lat 8. Casper Lat 6. Sobek Lat 3. \\
Marysia Rok 1.}
\end{centeredquote}

\begin{figure}[!ht]
    \centering \includegraphics[width=1.0\linewidth]{
        1760_Paweł_Piotrkowicz_akt_chrztu_parafia_Mińsk_Mazowiecki_wpis_12.jpg}
    \captionsetup{format=hang}
    \caption{Akt chrztu Pawła Piotrkowicz -- par. Mińsk Mazowiecki
    1760~rok (12/1760) \cite{par_minsk0}.}
    \label{fig:ppiotrkowicz_1760}
\end{figure}

\begin{figure}[!ht]
    \centering \includegraphics[width=0.9\linewidth]{
        1762_Kacper_Piotrkowicz_akt_chrztu_parafia_Mińsk_Mazowiecki_wpis_13.jpg}
    \captionsetup{format=hang}
    \caption{Akt chrztu Kacpra Piotrkowicz -- par. Mińsk Mazowiecki 1762~rok
    (13/1762) \cite{par_minsk0}.}
    \label{fig:kpriotrkowicz_1762}
\end{figure}

\begin{figure}[!ht]
    \centering \includegraphics[width=0.9\linewidth]{
        1768_Sebastian_Piotrkowicz_akt_chrztu_parafia_Mińsk_Mazowiecki_wpis_33.jpg}
    \captionsetup{format=hang}
    \caption{Akt chrztu Sebastiana Piotrkowicza -- par. Mińsk Mazowiecki
    1768~rok (33/1768) \cite{par_minsk0}.}
    \label{fig:spiotrkowicz_1768}
\end{figure}

\begin{figure}[!ht]
    \vspace*{0.5cm}
    \centering \includegraphics[width=0.85\linewidth]{
        1752_Maciej_Okrzeja_akt_chrztu_parafia_Mińsk_Mazowiecki_wpis_75.jpg}
    \captionsetup{format=hang}
    \caption{Akt chrztu Macieja Okrzei -- par. Mińsk Mazowiecki 1752~rok (75/1752) 
    \cite{par_minsk0}.}
    \label{fig:mokrzeja_1752}
\end{figure}

\begin{figure}[!ht]
    \vspace*{1cm}
    \centering \includegraphics[width=1.0\linewidth]{
        1778_Maciej_Okrzeja_Jadwiga_Gańko_akt_ślubu_parafia_Mińsk_Mazowiecki_wpis_84.jpg}
    \captionsetup{format=hang}
    \caption{Akt ślubu Macieja Okrzei i~Jadwigi Gańko -- par. Mińsk Mazowiecki
    1778~rok (84/1778) 
    \cite{par_minsk0}.}
    \label{fig:mokrzeja_1778}
\end{figure}

Z~perspektywy historii rodziny Gomulskich z~Desna, kluczową postacią
w~rodzinie Andrzeja Okrzei jest jego syn
Maciej\indeksosoby{Maciej}{Okrzeja}{1752}{Dębe Wielkie} oraz jego potomkowie.
Maciej Okrzeja urodził się 20~lutego 1752~roku w~Dębem (patrz: ryc.
\ref{fig:mokrzeja_1752}), a~w~1778~roku ożenił się z~osiemnastoletnią Jadwigą
Gańko\indeksosoby[Gańko]{Jadwiga}{Okrzeja}{1760}{Kobierne} (ur.~1760~r.~--
zm.~?) z~pobliskiej wsi Kobierne (patrz: ryc. \ref{fig:mokrzeja_1778}). Para
dochowała się co~najmniej dziewiątki dzieci:

\begin{itemize}
\item Marianny\indeksosoby{Marianna}{Okrzeja}{1779}{Dębe Wielkie}
(ur.~1779~r.~-- zm.~1780~r.),
\item \textbf{Doroty}\indeksosoby[Okrzeja]{Dorota}{Piotrkowicz}{1781}{Dębe Wielkie}
(ur.~1781~r.~-- zm.~1856~r.) -- patrz: ryc. \ref{fig:dokrzeja_1781},
\item \textbf{Antoniego}\indeksosoby{Antoni}{Okrzeja}{1783}{Dębe Wielkie}
(ur.~1783~r.~-- zm.~1844~r.) -- patrz: ryc. \ref{fig:aokrzeja_1783},
\item Kazimierza\indeksosoby{Kazimierz}{Okrzeja}{1786}{Dębe Wielkie}
(ur.~1786~r.~-- zm.~1811~r.),
\item Heleny\indeksosoby{Helena}{Okrzeja}{1788}{Dębe Wielkie}
(ur.~1788~r.~-- zm.~?),
\item Szczepana\indeksosoby{Szczepan}{Okrzeja}{1790}{Dębe Wielkie}
(ur.~1790~r.~-- zm.~?),
\item Agnieszki\indeksosoby{Agnieszka}{Okrzeja}{1793}{Dębe Wielkie}
(ur.~1793~r.~-- zm.~1863~r.),
\item Marianny\indeksosoby{Marianna}{Okrzeja}{1794}{Dębe Wielkie}
(ur.~1794~r.~-- zm.~?),
\item Tomasza\indeksosoby{Tomasz}{Okrzeja}{1796}{Dębe Wielkie}
(ur.~1796~r.~-- zm.~1849~r.).
\end{itemize}

\begin{figure}[!ht]
    \centering \includegraphics[width=0.8\linewidth]{
        1781_Dorota_Okrzeja_Piotrkowicz_akt_chrztu_parafia_Mińsk_Mazowiecki_wpis_42.jpg}
    \captionsetup{format=hang}
    \caption{Akt chrztu Doroty Okrzei -- par. Mińsk Mazowiecki 1781~rok (42/1781) 
    \cite{par_minsk0}.}
    \label{fig:dokrzeja_1781}
\end{figure}

\begin{figure}[!ht]
    \vspace*{0.5cm}
    \centering \includegraphics[width=0.8\linewidth]{
        1783_Antoni_Okrzeja_akt_chrztu_parafia_Mińsk_Mazowiecki_wpis_52.jpg}
    \captionsetup{format=hang}
    \caption{Akt chrztu Antoniego Okrzei -- par. Mińsk Mazowiecki 1783~rok
    (52/1783) \cite{par_minsk0}.}
    \label{fig:aokrzeja_1783}
\end{figure}

Natomiast w~rodzinie Bartłomieja Piotrkowicza potomkowie każdego z~jego trzech
synów: Pawła\indeksosoby{Paweł}{Piotrkowicz}{1760}{Dębe Wielkie},
Kacpra\indeksosoby{Kacper}{Piotrkowicz}{1762}{Dębe Wielkie}
i~Sebastiana\indeksosoby{Sebastian}{Piotrkowicz}{1768}{Dębe Wielkie} byli
związani z~rodziną Gomulskich z~Desna. W~przypadku Pawła i~Kacpra byli oni
również blisko związani z~rodziną Okrzejów -- Paweł Piotrkowicz ożenił się
z~Dorotą Okrzeją\indeksosoby[Okrzeja]{Dorota}{Piotrkowicz}{1781}{Dębe Wielkie}
córką Macieja Okrzei\indeksosoby{Maciej}{Okrzeja}{1752}{Dębe Wielkie}, a~córka
Kacpra Piotrkowicza, Zofia\indeksosoby[Piotrkowicz]{Zofia}{Okrzeja}{1791}{Dębe Wielkie}
(ur.~1791~r.~-- zm.~1832~r.), ożeniła się z~Antonim
Okrzeją\indeksosoby{Antoni}{Okrzeja}{1783}{Dębe Wielkie}, synem Macieja
Okrzei\indeksosoby{Maciej}{Okrzeja}{1752}{Dębe Wielkie}. Możemy więc wyróżnić
trzy główne gałęzie rodzinnych powiązań pomiędzy rodzinami Piotrkowiczów,
Okrzejów i~Gomulskich z~Desna: 

\begin{itemize}
    \item potomkowie Pawła\indeksosoby{Paweł}{Piotrkowicz}{1760}{Dębe Wielkie}
    i~Doroty\indeksosoby[Okrzeja]{Dorota}{Piotrkowicz}{1781}{Dębe Wielkie}
    Piotrkowiczów,
    \item potomkowie Antoniego\indeksosoby{Antoni}{Okrzeja}{1783}{Dębe Wielkie}
    i~Zofii\indeksosoby[Piotrkowicz]{Zofia}{Okrzeja}{1791}{Dębe Wielkie} Okrzejów,
    \item potomkowie Sebastiana\indeksosoby{Sebastian}{Piotrkowicz}{1768}{Dębe Wielkie}
    i~Anny\indeksosoby[Smuga]{Anna}{Piotrkowicz}{1770}{Dębe Wielkie} Piotrkowiczów.
\end{itemize}

Właśnie przez pryzmat tych gałęzi zaprezentuję w~dalszej części rozdziału
wspólną historię tych trzech rodzin. 

\BigBold{Potomkowie Pawła i~Doroty Piotrkowiczów}
\phantomsection\label{sec:pd_piotrkowicze}

Dnia 14~lutego 1801~roku w~parafii pod wezwaniem Narodzenia Najświętszej Maryi
Panny w~Mińsku Mazowieckim odbył się ślub pomiędzy najstarszym synem
Bartłomieja Piotrkowicza,
Pawłem\indeksosoby{Paweł}{Piotrkowicz}{1760}{Dębe Wielkie}, a~drugą córką
Macieja i~Jadwigi Okrzejów,
Dorotą\indeksosoby[Okrzeja]{Dorota}{Piotrkowicz}{1781}{Dębe Wielkie} (patrz
ryc. \ref{fig:dokrzeja_1801}). 

\begin{figure}[!ht]
    \centering \includegraphics[width=0.8\linewidth]{
        1801_Paweł_Piotrkowicz_Dorota_Okrzeja_akt_ślubu_parafia_Mińsk_Mazowiecki_wpis_21.jpg}
    \captionsetup{format=hang}
    \caption{Wpis dotyczący ślubu Pawła Piotrkowicza oraz Doroty Okrzei
    z~1801~roku w~parafii pod wezwaniem Narodzenia Najświętszej Maryi Panny
    w~Mińsku Mazowieckim   \cite{par_minsk0}.}
    \label{fig:dokrzeja_1801}
\end{figure}

\begin{figure}[!ht]
    \centering \includegraphics[width=0.78\linewidth]{
        1829_Tomasz_Piotrkowicz_Marianna_Lisiewska_akt_ślubu_parafia_Mińsk_Mazowiecki_wpis_30.jpg}
    \captionsetup{format=hang}
    \caption{Akt ślubu Tomasz Piotrkowicza i~Marianny Lisiewskiej -- par. Mińsk Mazowiecki
    1829~rok (30/1829)
    \cite{par_minsk2}.}
    \label{fig:tpiotrkowicz_1829}
\end{figure}

\begin{figure}[!ht]
\vspace*{0.5cm}
\centering
\begin{tikzpicture}[
  % Definiujemy stałe odstępy
  node distance=1cm and 1cm, 
  every node/.style={draw, rectangle, align=center, minimum width=3.8cm,
  inner sep=5pt}]

% --- POKOLENIE 0 ---
\node (Osiowa0) [draw=none, minimum width=0pt, inner sep=0pt] at (0,0) {};
\node (Andrzej) [left=1.5cm of Osiowa0] {Andrzej Okrzeja \\ {
    \footnotesize \emph{Dębe, 1721 r.}}};
\node (Katarzyna) [right=1.5cm of Osiowa0] {Katarzyna Brzost \\ {
    \footnotesize \emph{Brzóze, 1719 r.}}};

\draw[-] (Andrzej) -- (Katarzyna)
    node[midway, below, draw=none, minimum width=0pt,
    font=\scriptsize] (Slub0) {\emph{Mińsk, 1746 r.}};

% --- POKOLENIE 1 ---
\node (Osiowa1) [draw=none, minimum width=0pt, inner sep=0pt,
below=1.5cm of Slub0] {};
\node (Maciej) [left=1.5cm of Osiowa1] {Maciej Okrzeja \\ {
    \footnotesize \emph{Dębe, 1752 r.}}};
\node (Jadwiga) [right=1.5cm of Osiowa1] {Jadwiga Gańko \\ {
    \footnotesize \emph{Kobierne, 1760 r.}}};

\draw[-] (Maciej) -- (Jadwiga)
    node[midway, above, draw=none, minimum width=0pt,
    font=\scriptsize] (Slub1_Gora) { }
    node[midway, below, draw=none, minimum width=0pt,
    font=\scriptsize] (Slub1) {\emph{Mińsk, 1778 r.}};

\draw[->] (Slub0) -- (Slub1_Gora);

% --- POKOLENIE 2 ---
\node (Osiowa2) [draw=none, minimum width=0pt, inner sep=0pt,
below=1.5cm of Slub1] {};
\node (Dorota) [left=1.5cm of Osiowa2] {Dorota Okrzeja \\ {
    \footnotesize \emph{Dębe, 1781 r.}}};
\node (Pawel) [right=1.5cm of Osiowa2] {Paweł Piotrkowicz \\ {
    \footnotesize \emph{Dębe, 1760 r.}}};
    
\draw[-] (Dorota) -- (Pawel)
    node[midway, above, draw=none, minimum width=0pt,
    font=\scriptsize] (Slub2_Gora) { }
    node[midway, below, draw=none, minimum width=0pt,
    font=\scriptsize] (Slub2) {\emph{Mińsk, 1801 r.}};

\draw[->] (Slub1) -- (Slub2_Gora);

% --- POKOLENIE 3 ---
\node (Osiowa3) [draw=none, minimum width=0pt, inner sep=0pt,
below=1.5cm of  Slub2] {};
\node (Tomasz) [left=1.5cm of Osiowa3] {Tomasz Piotrkowicz \\ {
    \footnotesize \emph{Dębe, 1804 r.}}};
\node (Marianna) [right=1.5cm of Osiowa3] {Marianna Lisiewska \\ {
    \footnotesize \emph{Dębe, 1811 r.}}};

\draw[-] (Tomasz) -- (Marianna)
    node[midway, above, draw=none, minimum width=0pt,
    font=\scriptsize] (Slub3_Gora) { }
    node[midway, below, draw=none, minimum width=0pt,
    font=\scriptsize] (Slub3) {\emph{Mińsk, 1829 r.}};

\draw[->] (Slub2) -- (Slub3_Gora);

% --- POKOLENIE 4 ---
\node (Osiowa4) [draw=none, minimum width=0pt, inner sep=0pt,
below=1.5cm of  Slub3] {};
\node (Agnieszka) [left=1.5cm of Osiowa4] {Agnieszka Piotrkowicz \\
{\footnotesize \emph{Budy Barcząckie, 1840 r.}}};
\node (Jan) [right=1.5cm of Osiowa4] {
    Jan Gomulski \\ {\footnotesize \emph{Zamienie, 1836 r.}}};

\draw[-] (Agnieszka) -- (Jan)
    node[midway, above, draw=none, minimum width=0pt,
    font=\scriptsize] (Slub4_Gora) { }
    node[midway, below, draw=none, minimum width=0pt,
    font=\scriptsize] (Slub4) {\emph{Mińsk, 1858 r.}};

\draw[->] (Slub3) -- (Slub4_Gora);

% --- POKOLENIE 5 ---
\node (Osiowa5) [draw=none, minimum width=0pt, inner sep=0pt,
below=1.5cm of Slub4] {};
\node (Stanislaw) [left=1.5cm of Osiowa5] {Stanisław Gomulski \\
{\footnotesize \emph{Wólka Mińska, 1861 r.}}};
\node (Julianna) [right=1.5cm of Osiowa5] {
    Julianna Gągol \\ {\footnotesize \emph{Dłużka, 1862 r.}}};

\draw[-] (Stanislaw) -- (Julianna)
    node[midway, above, draw=none, minimum width=0pt,
    font=\scriptsize] (Slub5_Gora) { }
    node[midway, below, draw=none, minimum width=0pt,
    font=\scriptsize] (Slub5) {\emph{Mińsk, 1881 r.}};

\draw[->] (Slub4) -- (Slub5_Gora);

% --- POKOLENIE 6 ---
\node (Osiowa6) [draw=none, minimum width=0pt, inner sep=0pt,
below=1.5cm of Slub5] {};
\node (Jozef) [left=2cm of Osiowa6] {Józef Gomulski \\
{\footnotesize \emph{Wólka Mińska, 1895 r.}}};
\node (Feliksa) [right=2cm of Osiowa6] {
    Feliksa Tkaczyk \\ {\footnotesize \emph{Desno, 1897 r.}}};

\draw[-] (Jozef) -- (Feliksa)
    node[midway, above, draw=none, minimum width=0pt,
    font=\scriptsize] (Slub6_Gora) { }
    node[midway, below, draw=none, minimum width=0pt,
    font=\scriptsize] (Slub6) {\emph{Długa Kościelna, 1916 r.}};

\draw[->] (Slub5) -- (Slub6_Gora);

% --- POKOLENIE 7 ---
\node (Osiowa7) [draw=none, minimum width=0pt, inner sep=0pt,
below=1.5cm of Slub6] {};
\node (Stanislaw2) [left=2cm of Osiowa7] {Stanisław Gomulski \\
{\footnotesize \emph{Desno, 1926 r.}}};
\node (Jadwiga) [right=2cm of Osiowa7] {
    Jadwiga Leokadia \\ Poziemska \\ {\footnotesize \emph{Desno, 1935 r.}}};

\draw[-] (Stanislaw2) -- (Jadwiga)
    node[midway, above, draw=none, minimum width=0pt,
    font=\scriptsize] (Slub7_Gora) { }
    node[midway, below, draw=none, minimum width=0pt,
    font=\scriptsize] (Slub7) {\emph{Długa Kościelna, 1951 r.}};

\draw[->] (Slub6) -- (Slub7_Gora);

\end{tikzpicture}
\end{figure}

Małżeństwo Piotrkowiczów dochowało się wspólnie co~najmniej pięciorga dzieci:
Tomasza\indeksosoby{Tomasz}{Piotrkowicz}{1804}{Dębe Wielkie} (ur.~1804~r.~--
zm.~1874~r.), Marianny\indeksosoby{Marianna}{Piotrkowicz}{1813}{Dębe Wielkie}
(ur.~1813~r.~-- ?), Jakuba\indeksosoby{Jakub}{Piotrkowicz}{1816}{Dębe Wielkie}
(ur.~1816~r.~-- zm.~?), Józefy\indeksosoby{Józefa}{Piotrkowicz}{1819}{Dębe Wielkie}
(ur.~1819~r.~-- zm.~?) oraz Katarzyny\indeksosoby{Katarzyna}{Piotrkowicz}{1820}{Dębe Wielkie}
(ur.~1820~r.~-- zm.~?). Najstarszy syn Pawła i~Doroty, Tomasz, dnia 1~marca
1829~roku ożenił się z~Marianną
Lisiewską\indeksosoby[Lisiewska]{Marianna}{Piotrkowicz}{1811}{Dębe Wielkie}
(ur.~1811~r.~-- zm.~1874~r.) osiemnastoletnią wówczas panną z~Dębego (patrz
ryc. \ref{fig:tpiotrkowicz_1829}). Para dochowała się jedynie dwójki dzieci:
Agnieszki\indeksosoby[Piotrkowicz]{Agnieszka}{Gomulska}{1840}{Budy Barcząckie}
(ur.~1840~r.~-- zm.~1881~r.) oraz
Wojciecha\indeksosoby{Wojciech}{Piotrkowicz}{1845}{Budy Barcząckie}
(ur.~1845~r.~-- zm.~?). Agnieszka Piotrkowicz to pierwsza żona
Jana Gomulskiego\indeksosoby{Jan}{Gomulski}{1836}{Zamienie} (ur.~1836~r.~--
zm.~1901~r.) z~Wólki Mińskiej, których historię opisywałem w~rozdziale
IV. Tomasz Piotrowicz\indeksosoby{Tomasz}{Piotrkowicz}{1804}{Dębe Wielkie},
jego żona Marianna\indeksosoby[Lisiewska]{Marianna}{Piotrkowicz}{1811}{Dębe Wielkie}
oraz ich córka Agnieszka\indeksosoby[Piotrkowicz]{Agnieszka}{Gomulska}{1840}{Budy Barcząckie}
są~przodkami najmłodszej gałęzi rodziny Gomulskich z~Desna.

\newpage
\BigBold{Potomkowie Antoniego i~Zofii Okrzejów}
\phantomsection\label{sec:kf_piotrkowicze}

Syn Macieja i~Jadwigi, Antoni
Okrzeja\indeksosoby{Antoni}{Okrzeja}{1783}{Dębe Wielkie}, w~lutym
1809~roku ożenił się z~Zofią
Piotrkowicz\indeksosoby[Piotrkowicz]{Zofia}{Okrzeja}{1791}{Dębe Wielkie}
(ur.~1791~r.~-- zm.~1832~r.) córką Kacpra
Piotrkowicza\indeksosoby{Kacper}{Piotrkowicz}{1762}{Dębe Wielkie}
i~Franciszki Wiercioch\indeksosoby[Wiercioch]{Franciszka}{Piotrkowicz}{1770}{Dębe Wielkie}
(ur.~1770~r.~-- zm.~1830~r.) -- patrz: ryc. \ref{fig:kpiotrkowicz_1788},
\ref{fig:zpiotrkowicz_1791} oraz \ref{fig:aokrzeja_1809}. Para
dochowała się wspólnie co~najmniej czwórki dzieci:

\begin{itemize}
\item Macieja\indeksosoby{Maciej}{Okrzeja}{1810}{Dębe Wielkie} (ur.~1810~r.~--
zm.~1849~r.),
\item Szymona\indeksosoby{Szymon}{Okrzeja}{1811}{Dębe Wielkie}
(ur.~1811~r.~-- zm.~?),
\item Tomasza \indeksosoby{Tomasz}{Okrzeja}{1814}{Dębe Wielkie}
(ur.~1814~r.~-- zm.~1862~r.),
\item \textbf{Jana}\indeksosoby{Jan}{Okrzeja}{1825}{Dębe Wielkie}
(ur.~1825~r.~-- zm.~1901~r.) -- patrz: ryc. \ref{fig:jokrzeja_1825}.
\end{itemize}

\begin{figure}[!ht]
    \vspace*{0.5cm}
    \centering \includegraphics[width=0.9\linewidth]{
        1788_Kacper_Piotrkowicz_Franciszka_Wiercioch_akt_ślubu_parafia_Mińsk_Mazowiecki_wpis_5.jpg}
    \captionsetup{format=hang}
    \caption{Akt ślubu Kacpra\protect\footnotemark Piotrkowicza i~Franciszki
    Wiercioch -- par. Mińsk Mazowiecki 1788~rok (5/1788) \cite{par_minsk2}}.
    \label{fig:kpiotrkowicz_1788}
\end{figure}
\footnotetext{W~akcie tym Kacprowi Piotrkowiczowi mylnie przypisano imię
Franciszek, prawdopodobnie myląc je~z~imieniem jego małżonki Franciszki
Wiercioch.}

\begin{figure}[!ht]
    \centering \includegraphics[width=0.9\linewidth]{
        1791_Zofia_Piotrkowicz_akt_chrztu_parafia_Mińsk_Mazowiecki_wpis_65.jpg}
    \captionsetup{format=hang}
    \caption{Akt chrztu Zofii Piotrkowicz -- par. Mińsk Mazowiecki 1791~rok
    (65/1791) \cite{par_minsk2}.}
    \label{fig:zpiotrkowicz_1791}
\end{figure}

\begin{figure}[!ht]
    \centering \includegraphics[width=1.0\linewidth]{
        1809_Antoni_Okrzeja_Zofia_Piotrkowicz_akt_ślubu_parafia_Mińsk_Mazowiecki_wpis_23.jpg}
    \captionsetup{format=hang}
    \caption{Wpis dotyczący ślubu Antoniego Okrzei i~Zofii Piotrkowicz
    z~1809~roku w~parafii pod wezwaniem Narodzenia Najświętszej Maryi Panny
    w~Mińsku Mazowieckim   \cite{par_minsk0}.}
    \label{fig:aokrzeja_1809}
\end{figure}

\begin{figure}[!ht]
    \vspace*{0.5cm}
    \centering \includegraphics[width=0.9\linewidth]{
        1825_Jan_Okrzeja_akt_chrztu_parafia_Mińsk_Mazowiecki_wpis_158.jpg}
    \captionsetup{format=hang}
    \caption{Akt chrztu Jana Okrzei -- par. Mińsk Mazowiecki 1825~rok (158/1825) 
    \cite{par_minsk1}.}
    \label{fig:jokrzeja_1825}
\end{figure}

Z~perspektywy historii rodziny Gomulskich z~Desna, kluczową postacią
w~rodzinie Antoniego Okrzei\indeksosoby{Antoni}{Okrzeja}{1783}{Dębe Wielkie}
jest jego najmłodszy syn Jan\indeksosoby{Jan}{Okrzeja}{1825}{Dębe Wielkie},
który urodził się 26~maja 1825~roku w~Dębem (patrz: ryc.
\ref{fig:jokrzeja_1825}). Jan Okrzeja  dnia 17~listopada 1844~roku ożenił się
z~Marianną Gańko\indeksosoby[Gańko]{Marianna}{Okrzeja}{1824}{Dębe Wielkie}
(ur.~1824~r.~-- zm.~1854~r.) dwudziestoletnią panną z~jego rodzinnej wsi --
patrz: ryc. \ref{fig:jokrzeja_1844}.

\begin{figure}[!ht]
    \centering \includegraphics[width=0.8\linewidth]{
        1844_Jan_Okrzeja_Marianna_Gańko_akt_ślubu_parafia_Mińsk_Mazowiecki_wpis_41.jpg}
    \captionsetup{format=hang}
    \caption{Akt ślubu Jana Okrzei i~Marianny Gańko -- par. Mińsk Mazowiecki
    1844~rok (41/1844) 
    \cite{par_minsk2}.}
    \label{fig:jokrzeja_1844}
\end{figure}

\begin{figure}[!ht]
    \vspace*{0.5cm}
    \centering \includegraphics[width=1.0\linewidth]{
        1854_Marianna_Okrzeja_Bąk_akt_zgonu_parafia_Wiązowna_wpis_27.jpg}
    \captionsetup{format=hang}
    \caption{Akt zgonu Marianny Okrzei -- par. Wiązowna 1854~rok (27/1854)
    \cite{par_wiazowna}.}
    \label{fig:mokrzeja_1854}
\end{figure}

Jan\indeksosoby{Jan}{Okrzeja}{1825}{Dębe Wielkie} i~Marianna
Okrzejowie\indeksosoby[Gańko]{Marianna}{Okrzeja}{1824}{Dębe Wielkie} dochowali
się jedynie trójki potomstwa: Anny\indeksosoby{Anna}{Okrzeja}{1847}{Dębe Wielkie}
(ur.~1847~r.~-- zm.~?),
Magdaleny\indeksosoby{Magdalena}{Okrzeja}{1847}{Dębe Wielkie} (ur.~1850~r.~--
zm.~?) oraz Stanisława\indeksosoby{Stanisław}{Okrzeja}{1853}{Dębe Wielkie}
(ur.~1853~r.~-- zm.~1884~r.), gdyż 20~marca 1854~roku, po~zaledwie
10~latach ich małżeństwa,
Marianna\indeksosoby[Gańko]{Marianna}{Okrzeja}{1824}{Dębe Wielkie} zmarła
(patrz: ryc. \ref{fig:mokrzeja_1854}). Zgon Marianny Okrzei został odnotowany
w~księgach metrykalnych parafii pod wezwaniem św.~Wojciecha biskupa
i~męczennika w~Wiązownie -- jako miejsce jej śmierci wskazano wieś Konik
(obecnie Konik Stary) oddaloną o~około 8~kilometrów na zachód od~ich rodzinnej
wsi Dębe. Zmiana miejsca zamieszkania małżeństwa Okrzejów była prawdopodobnie
spowodowana faktem, iż~Jan Okrzeja około 1854~roku został dróżnikiem drogowym
trzeciej klasy dla Traktu Brzeskiego -- pierwszej bitej drogi w~Królestwie
Polskim, która przebiegała zarówno przez Dębe, jak i~Konik. Wskazuje
na~to~fakt, iż trójka najstarszych dzieci Jana urodziła się w~Dębem,
a~kolejne jego dzieci, zrodzone już z~nowego małżeństwa, przyszły na~świat
w~sześciu różnych miejscowościach zlokalizowanych wzdłuż Traktu Brzeskiego.
Dodatkowo o~nowym zawodzie\footnote{Jeszcze w~1853~roku, w~akcie chrztu
swojego syna Stanisława\indeksosoby{Stanisław}{Okrzeja}{1853}{Dębe Wielkie}
(tak jak w~poprzednich aktach metrykalnych) Jan
Okrzeja\indeksosoby{Jan}{Okrzeja}{1825}{Dębe Wielkie} określany jest jako
gospodarz / wyrobnik z~Dębego.} Jana
Okrzei\indeksosoby{Jan}{Okrzeja}{1825}{Dębe Wielkie} informują wprost zapisy
aktu jego drugiego ślubu, z~dnia 10~lipca 1854~roku, z~Marianną
Bok\indeksosoby[Bok]{Marianna}{Okrzeja}{1835}{Brzeziny}\footnote{Co~ciekawe
Marianna była starszą siostrą Franciszka Boka jednego ze~współzałożycieli
spółki Desna.} (ur.~1835~r.~-- zm.~1911~r.) panną mieszkającą w~sąsiedniej wsi
Brzeziny~-- patrz: ryc. \ref{fig:jokrzeja_1854}:

\begin{centeredquote}
\textit{... religijne małżeństwo między Janem Okrzeją, wdowcem po niegdyś
Mariannie z~Gańków Okrzejowej (...) dróżnikiem klasy trzeciej zamieszkałym
poprzednio w~Dembem Wielkim a~obecnie w~Koniku Duchnowskim a panną Marianną
Bok ...}
\end{centeredquote}

Jan Okrzeja\indeksosoby{Jan}{Okrzeja}{1825}{Dębe Wielkie} wziął więc drugi
ślub po~niecałych czterech miesiącach od~śmierci swojej pierwszej żony,
również Marianny. Z~drugą żoną dochował się co~najmniej dwanaściorga dzieci:
Marianny\indeksosoby{Marianna}{Okrzeja}{1855}{Konik Stary} (ur.~1855~r.~--
zm.~1878~r.), Franciszka\indeksosoby{Franciszek}{Okrzeja}{1857}{Konik Stary}
(ur.~1857~r.~-- zm.~1911~r.),
Walentego\indeksosoby{Walenty}{Okrzeja}{1860}{Janówek} (ur.~1860~r.~--
zm.~1914~r.), Joachima\indeksosoby{Joachim}{Okrzeja}{1862}{Janówek}
(ur.~1862~r.~-- zm.~?), Jana\indeksosoby{Jan}{Okrzeja}{1864}{Janówek}
(ur.~1864~r.~-- zm.~1864~r.), Józefa\indeksosoby{Józef}{Okrzeja}{1865}{Janówek}
(ur.~1865~r.~-- zm.~?), Rozalii\indeksosoby{Rozalia}{Okrzeja}{1868}{Janówek}
(ur.~1868~r.~-- zm.~?), Piotra\indeksosoby{Piotr}{Okrzeja}{1870}{Żurawka}
(ur.~1870~r.~-- zm.~?), Aleksandra\indeksosoby{Aleksander}{Okrzeja}{1873}{Żurawka}
(ur.~1873~r.~-- zm.~?), Michała\indeksosoby{Michał}{Okrzeja}{1870}{Miłosna}
(ur.~1875~r.~-- zm.~1908~r.), Stanisława\indeksosoby{Stanisław}{Okrzeja}{1878}{Hipolitów}
(ur.~1878~r.~-- zm.~?) oraz Antoniego
Władysława\indeksosoby{Antoni Władysław}{Okrzeja}{1884}{Hipolitów}
(ur.~1884~r.~-- zm.~?).

\begin{figure}[!ht]
    \centering \includegraphics[width=0.77\linewidth]{
        1854_Jan_Okrzeja_Marianna_Bok_akt_ślubu_parafia_Długa_Kościelna_wpis_4.jpg}
    \captionsetup{format=hang}
    \caption{Akt ślubu Jana Okrzei i~Marianny Bok -- par. Długa Kościelna
    1854~rok (4/1854) \cite{par_dluga}.}
    \label{fig:jokrzeja_1854}
\end{figure}

\begin{figure}[!ht]
    \vspace*{0.5cm}
    \centering \includegraphics[width=1.0\linewidth]{
        1901_Jan_Okrzeja_akt_zgonu_parafia_Długa_Kościelna_wpis_53.jpg}
    \captionsetup{format=hang}
    \caption{Akt zgonu Jana Okrzei -- par. Długa Kościelna 1901~rok
    (53/1901) \cite{par_dluga}.}
    \label{fig:jokrzeja_1901}
\end{figure}

\begin{figure}[!ht]
    \vspace*{0.5cm}
    \centering \includegraphics[width=1.0\linewidth]{
        1911_Marianna_Okrzeja_akt_zgonu_parafia_Długa_Kościelna_wpis_39.jpg}
    \captionsetup{format=hang}
    \caption{Akt zgonu Marianny Okrzei -- par. Długa Kościelna 1911~rok
    (39/1911) \cite{par_dluga}.}
    \label{fig:mokrzeja_1911}
\end{figure}

Ze względu na charakter wykonywanego zawodu przez Jana
Okrzeję\indeksosoby{Jan}{Okrzeja}{1825}{Dębe Wielkie}, czyli dbanie
o~utrzymanie porządku na~Trakcie Brzeskim, jego piętnaścioro dzieci z~obu
małżeństw urodziło się w~sześciu różnych miejscowościach. Jak już
wspominałem, trójka najstarszych dzieci, z~pierwszego małżeństwa, przyszła na
świat w~rodzinnej wsi Jana -- w~Dębem.
Marianna\indeksosoby{Marianna}{Okrzeja}{1855}{Konik Stary}
i~Franciszek\indeksosoby{Franciszek}{Okrzeja}{1857}{Konik Stary} urodzili się
w~Koniku, Walenty\indeksosoby{Walenty}{Okrzeja}{1860}{Janówek},
Joachim\indeksosoby{Joachim}{Okrzeja}{1862}{Janówek},
Jan\indeksosoby{Jan}{Okrzeja}{1864}{Janówek},
Józef\indeksosoby{Józef}{Okrzeja}{1865}{Janówek} oraz
Rozalia\indeksosoby{Rozalia}{Okrzeja}{1868}{Janówek} przyszli na~świat
w~Janówku (część obecnej wsi Zakręt).
Piotr\indeksosoby{Piotr}{Okrzeja}{1870}{Żurawka}
i~Aleksander\indeksosoby{Aleksander}{Okrzeja}{1873}{Żurawka} urodzili się
w~Żurawce (część obecnego Sulejówka),
Michał\indeksosoby{Michał}{Okrzeja}{1870}{Miłosna} w~Miłośnie (obecnie Stara
Miłosna), a~dwaj najmłodsi synowie Jana Okrzei,
Stanisław\indeksosoby{Stanisław}{Okrzeja}{1878}{Hipolitów} oraz~Antoni
Władysław\indeksosoby{Antoni Władysław}{Okrzeja}{1884}{Hipolitów}, przyszli
na~świat w~Hipolitowie. I~to właśnie w~Hipolitowie, około 1878~roku,
Jan Okrzeja\indeksosoby{Jan}{Okrzeja}{1825}{Dębe Wielkie} wraz ze~swoją liczną
rodziną osiedlił się na~stałe -- w~akcie jego zgonu z~1901~roku (patrz: ryc.
\ref{fig:jokrzeja_1901}) oraz zgonu jego drugiej żony z~1911~roku (patrz: ryc.
\ref{fig:mokrzeja_1911}) wpisana jest właśnie ta~miejscowość.

\begin{figure}[!ht]
    \vspace*{0.5cm}
    \centering \includegraphics[width=1.0\linewidth]{
        Ballada_ludowa_o_Stefanie_Okrzei.png}
    \captionsetup{format=hang}
    \caption{Ballada ludowa o~Stefanie Okrzei \cite{ballada_okrzeja}.}
    \label{fig:okrzeja_ballada}
\end{figure}

Spośród piętnaściorga dzieci Jana
Okrzei\indeksosoby{Jan}{Okrzeja}{1825}{Dębe Wielkie}, najbardziej interesujące
wydają losy potomków jego dwóch synów:
Franciszka\indeksosoby{Franciszek}{Okrzeja}{1857}{Konik Stary} oraz
Walentego\indeksosoby{Walenty}{Okrzeja}{1860}{Janówek}. Wnuk starszego
z~nich, Franciszka, Jan Ferdynand
Olszewski\indeksosoby{Jan Ferdynand \flagaPL}{Olszewski}{1930}{Warszawa}
(ur. 1930~r.~-- zm. 2019~r.) pełnił w~latach 1991--1992 funkcję Prezesa Rady
Ministrów Rzeczypospolitej Polskiej, był wybitnym prawnikiem, adwokatem oraz
działaczem opozycyjnym w~czasach Polskiej Rzeczpospolitej Ludowej. Natomiast
najstarszy syn Walentego Okrzei, Stefan Aleksander
Okrzeja\indeksosoby{Stefan Aleksander \flagaPL}{Okrzeja}{1886}{Dębe Wielkie}
(ur.~1884~r.~-- zm. 1905~r.), był bohaterem narodowym, działaczem
niepodległościowym i~socjalistycznym, który stał się symbolem rewolucji
robotniczej 1905~roku, a~jego męczeńska śmierć na stokach Cytadeli
Warszawskiej została upamiętniona w~wielu utworach literackich m.in.
w~\emph{Balladzie ludowej o~Stefanie Okrzei} (patrz: ryc.
\ref{fig:okrzeja_ballada}).

\begin{figure}[!ht]
    \vspace*{0.5cm}
    \centering \includegraphics[width=0.93\linewidth]{
        1857_Franciszek_Okrzeja_akt_chrztu_parafia_Wiązowna_wpis_19.jpg}
    \captionsetup{format=hang}
    \caption{Akt chrztu Franciszka Okrzei -- par. Wiązowna 1857~rok (19/1857)
    \cite{par_wiazowna}.}
    \label{fig:fokrzeja_1857}
\end{figure}

\begin{figure}[!ht]
    \vspace*{0.5cm}
    \centering \includegraphics[width=0.93\linewidth]{
        1860_Walenty_Okrzeja_akt_chrztu_parafia_Wiązowna_wpis_21.jpg}
    \captionsetup{format=hang}
    \caption{Akt chrztu Walentego Okrzei -- par. Wiązowna 1860~rok (21/1860)
    \cite{par_wiazowna}.}
    \label{fig:wokrzeja_1860}
\end{figure}

Franciszek Okrzeja\indeksosoby{Franciszek}{Okrzeja}{1857}{Konik Stary} dnia
21~września 1884~roku ożenił się z~Wiktorią
Jastrzębską\indeksosoby[Jastrzębska]{Wiktoria}{Okrzeja}{1862}{Warszawa}
(ur.~1862~r.~-- zm.~1931~r.) z~Warszawy (patrz: ryc. \ref{fig:fokrzeja_1884}),
para zamieszkała po~ślubie w~osadzie Bródno. Z~ich związku narodziło się
co~najmniej siedmioro dzieci: Helena
Eugenia\indeksosoby{Helena Eugenia}{Okrzeja}{1885}{Warszawa}
(ur.~1885~r.~-- zm.~1987~r.),
Zofia\indeksosoby{Zofia}{Okrzeja}{1887}{Warszawa} (ur.~1887~r.~-- zm.~1951~r.),
Marianna\indeksosoby{Marianna}{Okrzeja}{1892}{Warszawa} (ur.~1892~r.~-- zm.~?),
Bolesław\indeksosoby{Bolesław}{Okrzeja}{1894}{Warszawa} (ur.~1894~r.~--
zm.~1976~r.), Józef\indeksosoby{Józef}{Okrzeja}{1897}{Warszawa} (ur.~1897~r.~--
zm.~?), Jadwiga\indeksosoby[Okrzeja]{Jadwiga}{Olszewska}{1898}{Warszawa}
(ur.~1898~r.~-- zm.~1978~r.) oraz Feliksa\indeksosoby{Feliksa}{Okrzeja}{1901}{Warszawa}
(ur.~1901~r.~-- zm.~1903~r.). Z~aktu chrztu drugiej córki małżonków
Okrzejów wynika, iż w~1887~roku Franciszek był pracownikiem fizycznym przy
Drodze Żelaznej Warszawsko-Wiedeńskiej, pierwszej w~dziejach Królestwa
Polskiego linii kolejowej, która została uruchomiona w~1845~roku. Franciszek
kontynuował więc infrastrukturalną karierę swojego ojca
Jana\indeksosoby{Jan}{Okrzeja}{1825}{Dębe Wielkie}.

\begin{figure}[!ht]
    \centering \includegraphics[width=0.78\linewidth]{
        1884_Franciszek_Okrzeja_Wiktoria_Jastrzębska_akt_ślubu_parafia_MB_Loretańskiej_Warszawa_wpis_204.jpg}
    \captionsetup{format=hang}
    \caption{Akt ślubu Franciszka Okrzei i~Wiktorii Jastrzębskiej -- par. 
    Matki Boskiej Loretańskiej w~Warszawie 1884~rok (204/1884) 
    \cite{par_mb_warszawa}.}
    \label{fig:fokrzeja_1884}
\end{figure}

\begin{figure}[!ht]
    \vspace*{0.5cm}
    \centering \includegraphics[width=0.75\linewidth]{
        1898_Jadwiga_Okrzeja_Olszewska_akt_chrztu_parafia_MB_Loretańskiej_Warszawa_wpis_1966.jpg}
    \captionsetup{format=hang}
    \caption{Akt chrztu Jadwigi Okrzei -- par. Matki Boskiej Loretańskiej
    w~Warszawie 1898~rok (1966/1898) \cite{par_mb_warszawa}.}
    \label{fig:jokrzeja_1898}
\end{figure}

Jadwiga Okrzeja\indeksosoby[Okrzeja]{Jadwiga}{Olszewska}{1898}{Warszawa},
córka Franciszka\indeksosoby{Franciszek}{Okrzeja}{1857}{Konik Stary}
i~Wiktorii\indeksosoby[Jastrzębska]{Wiktoria}{Okrzeja}{1862}{Warszawa}, dnia
26~grudnia 1922~roku w~parafii pod wezwaniem Matki Boskiej Różańcowej
w~Warszawie, wyszła za~mąż za~Ferdynanda
Olszewskiego\indeksosoby{Ferdynand}{Olszewski}{1887}{Ossów} (ur.~1887~r.~--
zm.~1976~r.) urodzonym w~Ossowie pod Warszawą z~Seweryna i~Anny (z~Putków).
Ferdynand w~akcie ich ślubu jest nazwany kolejarzem, Jadwiga również
pochodziła z~kolejarskiej rodziny. Para dochowała się jedynie dwójki dzieci:
Marii\indeksosoby[Olszewska]{Maria}{Janocha}{1923}{Warszawa}\footnote{Syn
Marii Michał Janocha jest od~2015~roku biskupem pomocniczym w~archidiecezji
warszawskiej.} (ur.~1923~r.~-- zm.~2021~r.) oraz Jana
Ferdynanda\indeksosoby{Jan Ferdynand \flagaPL}{Olszewski}{1930}{Warszawa}
(ur.~1930~r.~-- zm.~2019~r.).

\begin{figure}[!ht]
    \centering \includegraphics[width=0.77\linewidth]{
        1922_Jadwiga_Okrzeja_Ferdynand_Olszewski_akt_ślubu_parafia_MB_Różańcowej_Warszawa_wpis_144.jpg}
    \captionsetup{format=hang}
    \caption{Akt ślubu Jadwigi Okrzei i~Ferdynanda Olszewskiego -- par. 
    Matki Boskiej Różańcowej w~Warszawie 1922~rok (144/1922) 
    \cite{par_mb2_warszawa}.}
    \label{fig:jokrzeja_1922}
\end{figure}

\begin{figure}[!ht]
\vspace*{0.5cm}
\centering
\begin{tikzpicture}[
  % Definiujemy stałe odstępy
  node distance=1cm and 1cm, 
  every node/.style={draw, rectangle, align=center, minimum width=3.8cm,
  inner sep=5pt}]

  % --- POKOLENIE 0 ---
\node (Osiowa0) [draw=none, minimum width=0pt, inner sep=0pt] at (0,0) {};
\node (Andrzej) [left=1.5cm of Osiowa0] {Andrzej Okrzeja \\ {
    \footnotesize \emph{Dębe, 1721 r.}}};
\node (Katarzyna) [right=1.5cm of Osiowa0] {Katarzyna Brzost \\ {
    \footnotesize \emph{Brzóze, 1719 r.}}};

\draw[-] (Andrzej) -- (Katarzyna)
    node[midway, below, draw=none, minimum width=0pt,
    font=\scriptsize] (Slub0) {\emph{Mińsk, 1746 r.}};

% --- POKOLENIE 1 ---
\node (Osiowa1) [draw=none, minimum width=0pt, inner sep=0pt,
below=1.5cm of Slub0] {};
\node (Maciej) [left=1.5cm of Osiowa1] {Maciej Okrzeja \\ {
    \footnotesize \emph{Dębe, 1752 r.}}};
\node (Jadwiga) [right=1.5cm of Osiowa1] {Jadwiga Gańko \\ {
    \footnotesize \emph{Kobierne, 1760 r.}}};

\draw[-] (Maciej) -- (Jadwiga)
    node[midway, above, draw=none, minimum width=0pt,
    font=\scriptsize] (Slub1_Gora) { }
    node[midway, below, draw=none, minimum width=0pt,
    font=\scriptsize] (Slub1) {\emph{Mińsk, 1778 r.}};

\draw[->] (Slub0) -- (Slub1_Gora);

% STRZAŁKA W DÓŁ
\draw[->] (Slub1.south) -- ++(0,-1.5);

\end{tikzpicture}
\end{figure}

\begin{figure}[!ht]
\centering
% --- PIERWSZY WYKRES ---
\begin{minipage}[t]{0.48\textwidth}
\centering
\begin{tikzpicture}[baseline=(Osiowa2.north), scale=0.7, transform shape,
    node distance=1cm and 1cm, 
    every node/.style={draw, rectangle, align=center, minimum width=3.5cm,
    inner sep=4pt}]

% --- POKOLENIE 2 ---
\node (Osiowa2) [draw=none, minimum width=0pt, inner sep=0pt] at (0,0) {};
\node (Dorota) [left=1.5cm of Osiowa2] {Dorota Okrzeja \\ {
    \footnotesize \emph{Dębe, 1781 r.}}};
\node (Pawel) [right=1.5cm of Osiowa2] {Paweł Piotrkowicz \\ {
    \footnotesize \emph{Dębe, 1760 r.}}};
    
\draw[-] (Dorota) -- (Pawel)
    node[midway, above, draw=none, minimum width=0pt,
    font=\scriptsize] (Slub2_Gora) { }
    node[midway, below, draw=none, minimum width=0pt,
    font=\scriptsize] (Slub2) {\emph{Mińsk, 1801 r.}};

% --- POKOLENIE 3 ---
\node (Osiowa3) [draw=none, minimum width=0pt, inner sep=0pt,
below=1.5cm of  Slub2] {};
\node (Tomasz) [left=1.5cm of Osiowa3] {Tomasz Piotrkowicz \\ {
    \footnotesize \emph{Dębe, 1804 r.}}};
\node (Marianna) [right=1.5cm of Osiowa3] {Marianna Lisiewska \\ {
    \footnotesize \emph{Dębe, 1811 r.}}};

\draw[-] (Tomasz) -- (Marianna)
    node[midway, above, draw=none, minimum width=0pt,
    font=\scriptsize] (Slub3_Gora) { }
    node[midway, below, draw=none, minimum width=0pt,
    font=\scriptsize] (Slub3) {\emph{Mińsk, 1829 r.}};

\draw[->] (Slub2) -- (Slub3_Gora);

% --- POKOLENIE 4 ---
\node (Osiowa4) [draw=none, minimum width=0pt, inner sep=0pt,
below=1.5cm of  Slub3] {};
\node (Agnieszka) [left=1.5cm of Osiowa4] {Agnieszka Piotrkowicz \\
{\footnotesize \emph{Budy Barcząckie, 1840 r.}}};
\node (Jan) [right=1.5cm of Osiowa4] {
    Jan Gomulski \\ {\footnotesize \emph{Zamienie, 1836 r.}}};

\draw[-] (Agnieszka) -- (Jan)
    node[midway, above, draw=none, minimum width=0pt,
    font=\scriptsize] (Slub4_Gora) { }
    node[midway, below, draw=none, minimum width=0pt,
    font=\scriptsize] (Slub4) {\emph{Mińsk, 1858 r.}};

\draw[->] (Slub3) -- (Slub4_Gora);

% --- POKOLENIE 5 ---
\node (Osiowa5) [draw=none, minimum width=0pt, inner sep=0pt,
below=1.5cm of Slub4] {};
\node (Stanislaw) [left=1.5cm of Osiowa5] {Stanisław Gomulski \\
{\footnotesize \emph{Wólka Mińska, 1861 r.}}};
\node (Julianna) [right=1.5cm of Osiowa5] {
    Julianna Gągol \\ {\footnotesize \emph{Dłużka, 1862 r.}}};

\draw[-] (Stanislaw) -- (Julianna)
    node[midway, above, draw=none, minimum width=0pt,
    font=\scriptsize] (Slub5_Gora) { }
    node[midway, below, draw=none, minimum width=0pt,
    font=\scriptsize] (Slub5) {\emph{Mińsk, 1881 r.}};

\draw[->] (Slub4) -- (Slub5_Gora);

% --- POKOLENIE 6 ---
\node (Osiowa6) [draw=none, minimum width=0pt, inner sep=0pt,
below=1.5cm of Slub5] {};
\node (Jozef) [left=2cm of Osiowa6] {Józef Gomulski \\
{\footnotesize \emph{Wólka Mińska, 1895 r.}}};
\node (Feliksa) [right=2cm of Osiowa6] {
    Feliksa Tkaczyk \\ {\footnotesize \emph{Desno, 1897 r.}}};

\draw[-] (Jozef) -- (Feliksa)
    node[midway, above, draw=none, minimum width=0pt,
    font=\scriptsize] (Slub6_Gora) { }
    node[midway, below, draw=none, minimum width=0pt,
    font=\scriptsize] (Slub6) {\emph{Długa Kościelna, 1916 r.}};

\draw[->] (Slub5) -- (Slub6_Gora);

% --- POKOLENIE 7 ---
\node (Osiowa7) [draw=none, minimum width=0pt, inner sep=0pt,
below=1.5cm of Slub6] {};
\node (Stanislaw2) [left=2cm of Osiowa7] {Stanisław Gomulski \\
{\footnotesize \emph{Desno, 1926 r.}}};
\node (Jadwiga) [right=2cm of Osiowa7] {
    Jadwiga Leokadia \\ Poziemska \\ {\footnotesize \emph{Desno, 1935 r.}}};

\draw[-] (Stanislaw2) -- (Jadwiga)
    node[midway, above, draw=none, minimum width=0pt,
    font=\scriptsize] (Slub7_Gora) { }
    node[midway, below, draw=none, minimum width=0pt,
    font=\scriptsize] (Slub7) {\emph{Długa Kościelna, 1951 r.}};

\draw[->] (Slub6) -- (Slub7_Gora);
   
\end{tikzpicture}
\end{minipage}
\hfill % Doda elastyczny odstęp między wykresami
% --- DRUGI WYKRES ---
\begin{minipage}[t]{0.48\textwidth}
\centering
\begin{tikzpicture}[baseline=(Osiowa2.north), scale=0.7, transform shape,
    node distance=1cm and 1cm, 
    every node/.style={draw, rectangle, align=center, minimum width=3.5cm,
    inner sep=4pt}]

% --- POKOLENIE 2 ---
\node (Osiowa2) [draw=none, minimum width=0pt, inner sep=0pt] at (0,0) {};
\node (Antoni) [left=1.5cm of Osiowa2] {Antoni Okrzeja \\ {
    \footnotesize \emph{Dębe, 1783 r.}}};
\node (Zofia) [right=1.5cm of Osiowa2] {Zofia Piotrkowicz \\ {
    \footnotesize \emph{Dębe, 1791 r.}}};
    
\draw[-] (Antoni) -- (Zofia)
    node[midway, above, draw=none, minimum width=0pt,
    font=\scriptsize] (Slub2_Gora) { }
    node[midway, below, draw=none, minimum width=0pt,
    font=\scriptsize] (Slub2) {\emph{Mińsk, 1809 r.}};

% --- POKOLENIE 3 ---
\node (Osiowa3) [draw=none, minimum width=0pt, inner sep=0pt,
below=1.5cm of  Slub2] {};
\node (Jan) [left=1.5cm of Osiowa3] {Jan Okrzeja \\ {
    \footnotesize \emph{Dębe, 1825 r.}}};
\node (Marianna) [right=1.5cm of Osiowa3] {Marianna Bok \\ {
    \footnotesize \emph{Jabłonna, 1835 r.}}};

\draw[-] (Jan) -- (Marianna)
    node[midway, above, draw=none, minimum width=0pt,
    font=\scriptsize] (Slub3_Gora) { }
    node[midway, below, draw=none, minimum width=0pt,
    font=\scriptsize] (Slub3) {\emph{Długa Kościelna, 1854 r.}};

\draw[->] (Slub2) -- (Slub3_Gora);

% --- POKOLENIE 4 ---
\node (Osiowa4) [draw=none, minimum width=0pt, inner sep=0pt,
below=1.5cm of  Slub3] {};
\node (Franciszek) [left=1.5cm of Osiowa4] {Franciszek Okrzeja \\
{\footnotesize \emph{Konik, 1857 r.}}};
\node (Wiktoria) [right=1.5cm of Osiowa4] {
    Wiktoria Jastrzębska \\ {\footnotesize \emph{Kowiesy, 1862 r.}}};

\draw[-] (Franciszek) -- (Wiktoria)
    node[midway, above, draw=none, minimum width=0pt,
    font=\scriptsize] (Slub4_Gora) { }
    node[midway, below, draw=none, minimum width=0pt,
    font=\scriptsize] (Slub4) {\emph{Warszawa, 1884 r.}};

\draw[->] (Slub3) -- (Slub4_Gora);

% --- POKOLENIE 5 ---
\node (Osiowa5) [draw=none, minimum width=0pt, inner sep=0pt,
below=1.5cm of Slub4] {};
\node (Jadwiga) [left=1.5cm of Osiowa5] {Jadwiga Okrzeja \\
{\footnotesize \emph{Warszawa, 1898 r.}}};
\node (Ferdynand) [right=1.5cm of Osiowa5] {
    Ferdynand Olszewski \\ {\footnotesize \emph{Ossów, 1886 r.}}};

\draw[-] (Jadwiga) -- (Ferdynand)
    node[midway, above, draw=none, minimum width=0pt,
    font=\scriptsize] (Slub5_Gora) { }
    node[midway, below, draw=none, minimum width=0pt,
    font=\scriptsize] (Slub5) {\emph{Warszawa, 1922 r.}};

\draw[->] (Slub4) -- (Slub5_Gora);

% --- POKOLENIE 6 ---
\node (JanFerdynand) [below=1.5cm of Slub5] {Jan Ferdynand Olszewski \\
{\footnotesize \emph{Warszawa, 1930 r.}} \\ PREZES RADY MINISTRÓW \flagaPL};

\draw[->] (Slub5) -- (JanFerdynand);

\end{tikzpicture}
\end{minipage}
\vspace*{0.5cm}
\end{figure}

Walenty Okrzeja\indeksosoby{Walenty}{Okrzeja}{1860}{Janówek} dnia 4~lutego
1885~roku ożenił się z~Helena
Cieślińską\indeksosoby[Cieślińska]{Helena}{Okrzeja}{1866}{Grodno}
(ur.~1866~r.~-- zm.~1949~r.) dziewiętnastoletnią wówczas panną z~Długiej
Kościelnej, urodzoną w~Grodnie, córką Karola i~Teofili (z~domu Dąbrowskiej) --
patrz: ryc. \ref{fig:wokrzeja_1885}. Ślub odbył się w~parafii św.~Anny
w~Długiej Kościelnej, do~której oboje młodzi należeli. Z~treści aktu ich ślubu
wynika, iż Walenty Okrzeja był wówczas pracownikiem kolejowym na nowo
powstałej trasie kolei Warszawsko-Terespolskiej\footnote{Trasa ta została
otwarta w~1867~roku.}, co było prawdopodobnie związane z~tym, iż przebiegała
ona przez Hipolitów, w~którym na stałe osiedliła się jego rodzina. Kontynuował
on więc tradycje zawodowe swojego ojca
Jana\indeksosoby{Jan}{Okrzeja}{1825}{Dębe Wielkie} oraz starszego brata
Franciszka\indeksosoby{Franciszek}{Okrzeja}{1857}{Konik Stary}.
Para początkowo zamieszkiwała w~rodzinnej wsi ojca Franciszka -- w~Dębem
(przez które również przebiegała trasa kolei Warszawsko-Terespolskiej), gdzie
urodziła się czwórka ich najstarszych dzieci, po czym około 1893~roku
przeprowadzili się do wsi Kamionek\footnote{Obecnie osiedle na~warszawskiej
Pradze.}, gdzie przyszło na świat kolejne ich potomstwo.
Walenty\indeksosoby{Walenty}{Okrzeja}{1860}{Janówek}
i~Helena\indeksosoby[Cieślińska]{Helena}{Okrzeja}{1866}{Grodno}
Okrzejowie dochowali się co~najmniej jedenaściorga dzieci:
Stefana Aleksandra\indeksosoby{Stefan Aleksander \flagaPL}{Okrzeja}{1886}{Dębe Wielkie}
(ur.~1886~r.~-- zm.~1905~r.),
Kazimiery Pelagii\indeksosoby[Okrzeja]{Kazimiera Pelagia}{Mieszczanek}{1888}{Dębe Wielkie}
(ur.~1888~r.~-- zm.~1976~r.),
Leonarda Józefa\indeksosoby{Leonard Józef}{Okrzeja}{1889}{Dębe Wielkie}
(ur.~1889~r.~-- zm.~1889~r.),
Wacława Zygmunta\indeksosoby{Wacław Zygmunt \flagaPL}{Okrzeja}{1891}{Dębe Wielkie}
(ur.~1891~r.~-- zm.~1955~r.),
Józefa Marcelego\indeksosoby{Józef Marceli}{Okrzeja}{1893}{Warszawa}
(ur.~1893~r.~-- zm.~?),
Marty\indeksosoby[Okrzeja]{Marta}{Czarniecka}{1895}{Warszawa}
(ur.~1895~r.~-- zm.~1976~r.),
Władysława\indeksosoby{Władysław \flagaPL}{Okrzeja}{1897}{Warszawa}
(ur.~1897~r.~-- zm.~1919~r.),
Michała Alfonsa\indeksosoby{Michał Alfons \flagaPL}{Okrzeja}{1899}{Warszawa}
(ur.~1899~r.~-- zm.~1942~r.),
Cecylii Elżbiety\indeksosoby{Cecylia Elżbieta}{Okrzeja}{1902}{Warszawa}
(ur.~1902~r.~-- zm.~1903~r.),
Jadwigi\indeksosoby[Okrzeja]{Jadwiga}{Strzelbiecka}{1904}{Warszawa}
(ur.~1904~r.~-- zm.~1996~r.) oraz
Stefana Stanisława\indeksosoby{Stefan Stanisław \flagaPL}{Okrzeja}{1907}{Warszawa}
(ur.~1907~r.~-- zm.~1939~r.).

\begin{figure}[!ht]
    \centering \includegraphics[width=0.87\linewidth]{
        1885_Walenty_Okrzeja_Helena_Cieślińska_akt_ślubu_parafia_Długa_Kościelna_wpis_4.jpg}
    \captionsetup{format=hang}
    \caption{Akt ślubu Walentego Okrzei i~Heleny Cieślińskiej -- par. 
    Długa Kościelna 1885~rok (4/1885) 
    \cite{par_dluga}.}
    \label{fig:wokrzeja_1885}
\end{figure}

\begin{figure}[!ht]
    \vspace*{1cm}
    \centering \includegraphics[width=1.0\linewidth]{
        1886_Stefan_Aleksander_Okrzeja_akt_chrztu_parafia_Mińsk_Mazowiecki_wpis_115.jpg}
    \captionsetup{format=hang}
    \caption{Akt chrztu Stefana Aleksandra Okrzei -- par. Mińsk Mazowiecki
    1886~rok (115/1886) \cite{par_minsk2}.}
    \label{fig:sokrzeja_1886}
\end{figure}

Nie tylko Stefan Aleksander
Okrzeja\indeksosoby{Stefan Aleksander \flagaPL}{Okrzeja}{1886}{Dębe Wielkie}
zasłużył się dla historii Polski, ale również inni członkowie jego
rodziny\footnote{Szczegółowa historia poszczególnych członków rodziny Okrzejów
została przedstawiona w~artykule prasowym Pawła Tomczyka opublikowanym
22~lipca 2025~roku na~portalu historycznym \emph{dzieje.pl} pod tytułem
\enquote{120~lat temu na~stokach warszawskiej Cytadeli stracono Stefana
Okrzeję} \cite{okrzeja_art}.}:

\begin{itemize}
\item Brat Władysław\indeksosoby{Władysław \flagaPL}{Okrzeja}{1897}{Warszawa}
zginął w~wieku 22~lat jako plutonowy podchorąży w~czasie wojny
polsko-bolszewickiej.
\item Brat Michał
Alfons\indeksosoby{Michał Alfons \flagaPL}{Okrzeja}{1899}{Warszawa}, jako
żołnierz walczył we~wrześniu 1939~roku, w~bitwie pod Garwolinem został ciężko
ranny~-- zmarł w~szpitalu w~1942~roku.
\item Kolejny z~braci Stefana, Wacław
Zygmunt\indeksosoby{Wacław Zygmunt \flagaPL}{Okrzeja}{1891}{Dębe Wielkie},
w~wieku 17~lat został aresztowany przez carską policję za~działalność
niepodległościową po~czym zesłano go~na~Syberię. Wrócił do~Polski przed
1920~rokiem, po~czym w~1944~roku uczestniczył w~Powstaniu Warszawskim pod
pseudonimem \enquote{Stefan} w~I~plutonie Zgrupowania Użyteczności Publicznej
WSOP \enquote{Elektrownia}. Po~upadku Powiśla pracował w~pionie elektrycznym
Delegatury Rządu w~Śródmieściu.
\item Najmłodszy brat Stefana, który dostał po~nim imię, Stefan
Stanisław\indeksosoby{Stefan Stanisław \flagaPL}{Okrzeja}{1907}{Warszawa}, był
porucznikiem pilotem Wojska Polskiego w~112~Eskadrze Myśliwskiej~-- zginął
w~powietrzu 5~września 1939~roku nad miejscowością Kręgi na~północnym Mazowszu
broniąc Warszawy przed niemieckim najeźdźcą.
\end{itemize}

\begin{figure}[!ht]
\vspace*{0.75cm}
\centering
\begin{tikzpicture}[
  % Definiujemy stałe odstępy
  node distance=1cm and 1cm, 
  every node/.style={draw, rectangle, align=center, minimum width=3.8cm,
  inner sep=5pt}]

  % --- POKOLENIE 0 ---
\node (Osiowa0) [draw=none, minimum width=0pt, inner sep=0pt] at (0,0) {};
\node (Andrzej) [left=1.5cm of Osiowa0] {Andrzej Okrzeja \\ {
    \footnotesize \emph{Dębe, 1721 r.}}};
\node (Katarzyna) [right=1.5cm of Osiowa0] {Katarzyna Brzost \\ {
    \footnotesize \emph{Brzóze, 1719 r.}}};

\draw[-] (Andrzej) -- (Katarzyna)
    node[midway, below, draw=none, minimum width=0pt,
    font=\scriptsize] (Slub0) {\emph{Mińsk, 1746 r.}};

% --- POKOLENIE 1 ---
\node (Osiowa1) [draw=none, minimum width=0pt, inner sep=0pt,
below=1.5cm of Slub0] {};
\node (Maciej) [left=1.5cm of Osiowa1] {Maciej Okrzeja \\ {
    \footnotesize \emph{Dębe, 1752 r.}}};
\node (Jadwiga) [right=1.5cm of Osiowa1] {Jadwiga Gańko \\ {
    \footnotesize \emph{Kobierne, 1760 r.}}};

\draw[-] (Maciej) -- (Jadwiga)
    node[midway, above, draw=none, minimum width=0pt,
    font=\scriptsize] (Slub1_Gora) { }
    node[midway, below, draw=none, minimum width=0pt,
    font=\scriptsize] (Slub1) {\emph{Mińsk, 1778 r.}};

\draw[->] (Slub0) -- (Slub1_Gora);

% STRZAŁKA W DÓŁ
\draw[->] (Slub1.south) -- ++(0,-1.5);

\end{tikzpicture}
\end{figure}

\begin{figure}[!ht]
\centering
% --- PIERWSZY WYKRES ---
\begin{minipage}[t]{0.48\textwidth}
\centering
\begin{tikzpicture}[baseline=(Osiowa2.north), scale=0.7, transform shape,
    node distance=1cm and 1cm, 
    every node/.style={draw, rectangle, align=center, minimum width=3.5cm,
    inner sep=4pt}]

% --- POKOLENIE 2 ---
\node (Osiowa2) [draw=none, minimum width=0pt, inner sep=0pt] at (0,0) {};
\node (Dorota) [left=1.5cm of Osiowa2] {Dorota Okrzeja \\ {
    \footnotesize \emph{Dębe, 1781 r.}}};
\node (Pawel) [right=1.5cm of Osiowa2] {Paweł Piotrkowicz \\ {
    \footnotesize \emph{Dębe, 1760 r.}}};
    
\draw[-] (Dorota) -- (Pawel)
    node[midway, above, draw=none, minimum width=0pt,
    font=\scriptsize] (Slub2_Gora) { }
    node[midway, below, draw=none, minimum width=0pt,
    font=\scriptsize] (Slub2) {\emph{Mińsk, 1801 r.}};

% --- POKOLENIE 3 ---
\node (Osiowa3) [draw=none, minimum width=0pt, inner sep=0pt,
below=1.5cm of  Slub2] {};
\node (Tomasz) [left=1.5cm of Osiowa3] {Tomasz Piotrkowicz \\ {
    \footnotesize \emph{Dębe, 1804 r.}}};
\node (Marianna) [right=1.5cm of Osiowa3] {Marianna Lisiewska \\ {
    \footnotesize \emph{Dębe, 1811 r.}}};

\draw[-] (Tomasz) -- (Marianna)
    node[midway, above, draw=none, minimum width=0pt,
    font=\scriptsize] (Slub3_Gora) { }
    node[midway, below, draw=none, minimum width=0pt,
    font=\scriptsize] (Slub3) {\emph{Mińsk, 1829 r.}};

\draw[->] (Slub2) -- (Slub3_Gora);

% --- POKOLENIE 4 ---
\node (Osiowa4) [draw=none, minimum width=0pt, inner sep=0pt,
below=1.5cm of  Slub3] {};
\node (Agnieszka) [left=1.5cm of Osiowa4] {Agnieszka Piotrkowicz \\
{\footnotesize \emph{Budy Barcząckie, 1840 r.}}};
\node (Jan) [right=1.5cm of Osiowa4] {
    Jan Gomulski \\ {\footnotesize \emph{Zamienie, 1836 r.}}};

\draw[-] (Agnieszka) -- (Jan)
    node[midway, above, draw=none, minimum width=0pt,
    font=\scriptsize] (Slub4_Gora) { }
    node[midway, below, draw=none, minimum width=0pt,
    font=\scriptsize] (Slub4) {\emph{Mińsk, 1858 r.}};

\draw[->] (Slub3) -- (Slub4_Gora);

% --- POKOLENIE 5 ---
\node (Osiowa5) [draw=none, minimum width=0pt, inner sep=0pt,
below=1.5cm of Slub4] {};
\node (Stanislaw) [left=1.5cm of Osiowa5] {Stanisław Gomulski \\
{\footnotesize \emph{Wólka Mińska, 1861 r.}}};
\node (Julianna) [right=1.5cm of Osiowa5] {
    Julianna Gągol \\ {\footnotesize \emph{Dłużka, 1862 r.}}};

\draw[-] (Stanislaw) -- (Julianna)
    node[midway, above, draw=none, minimum width=0pt,
    font=\scriptsize] (Slub5_Gora) { }
    node[midway, below, draw=none, minimum width=0pt,
    font=\scriptsize] (Slub5) {\emph{Mińsk, 1881 r.}};

\draw[->] (Slub4) -- (Slub5_Gora);

% --- POKOLENIE 6 ---
\node (Osiowa6) [draw=none, minimum width=0pt, inner sep=0pt,
below=1.5cm of Slub5] {};
\node (Jozef) [left=2cm of Osiowa6] {Józef Gomulski \\
{\footnotesize \emph{Wólka Mińska, 1895 r.}}};
\node (Feliksa) [right=2cm of Osiowa6] {
    Feliksa Tkaczyk \\ {\footnotesize \emph{Desno, 1897 r.}}};

\draw[-] (Jozef) -- (Feliksa)
    node[midway, above, draw=none, minimum width=0pt,
    font=\scriptsize] (Slub6_Gora) { }
    node[midway, below, draw=none, minimum width=0pt,
    font=\scriptsize] (Slub6) {\emph{Długa Kościelna, 1916 r.}};

\draw[->] (Slub5) -- (Slub6_Gora);

% --- POKOLENIE 7 ---
\node (Osiowa7) [draw=none, minimum width=0pt, inner sep=0pt,
below=1.5cm of Slub6] {};
\node (Stanislaw2) [left=2cm of Osiowa7] {Stanisław Gomulski \\
{\footnotesize \emph{Desno, 1926 r.}}};
\node (Jadwiga) [right=2cm of Osiowa7] {
    Jadwiga Leokadia \\ Poziemska \\ {\footnotesize \emph{Desno, 1935 r.}}};

\draw[-] (Stanislaw2) -- (Jadwiga)
    node[midway, above, draw=none, minimum width=0pt,
    font=\scriptsize] (Slub7_Gora) { }
    node[midway, below, draw=none, minimum width=0pt,
    font=\scriptsize] (Slub7) {\emph{Długa Kościelna, 1951 r.}};

\draw[->] (Slub6) -- (Slub7_Gora);
   
\end{tikzpicture}
\end{minipage}
\hfill % Doda elastyczny odstęp między wykresami
% --- DRUGI WYKRES ---
\begin{minipage}[t]{0.48\textwidth}
\centering
\begin{tikzpicture}[baseline=(Osiowa2.north), scale=0.7, transform shape,
    node distance=1cm and 1cm, 
    every node/.style={draw, rectangle, align=center, minimum width=3.5cm,
    inner sep=4pt}]

% --- POKOLENIE 2 ---
\node (Osiowa2) [draw=none, minimum width=0pt, inner sep=0pt] at (0,0) {};
\node (Antoni) [left=1.5cm of Osiowa2] {Antoni Okrzeja \\ {
    \footnotesize \emph{Dębe, 1783 r.}}};
\node (Zofia) [right=1.5cm of Osiowa2] {Zofia Piotrkowicz \\ {
    \footnotesize \emph{Dębe, 1791 r.}}};
    
\draw[-] (Antoni) -- (Zofia)
    node[midway, above, draw=none, minimum width=0pt,
    font=\scriptsize] (Slub2_Gora) { }
    node[midway, below, draw=none, minimum width=0pt,
    font=\scriptsize] (Slub2) {\emph{Mińsk, 1809 r.}};

% --- POKOLENIE 3 ---
\node (Osiowa3) [draw=none, minimum width=0pt, inner sep=0pt,
below=1.5cm of  Slub2] {};
\node (Jan) [left=1.5cm of Osiowa3] {Jan Okrzeja \\ {
    \footnotesize \emph{Dębe, 1825 r.}}};
\node (Marianna) [right=1.5cm of Osiowa3] {Marianna Bok \\ {
    \footnotesize \emph{Jabłonna, 1835 r.}}};

\draw[-] (Jan) -- (Marianna)
    node[midway, above, draw=none, minimum width=0pt,
    font=\scriptsize] (Slub3_Gora) { }
    node[midway, below, draw=none, minimum width=0pt,
    font=\scriptsize] (Slub3) {\emph{Długa Kościelna, 1854 r.}};

\draw[->] (Slub2) -- (Slub3_Gora);

% --- POKOLENIE 4 ---
\node (Osiowa4) [draw=none, minimum width=0pt, inner sep=0pt,
below=1.5cm of  Slub3] {};
\node (Walenty) [left=1.5cm of Osiowa4] {Walenty Okrzeja \\
{\footnotesize \emph{Janówek, 1860 r.}}};
\node (Helena) [right=1.5cm of Osiowa4] {
    Helena Cieślińska \\ {\footnotesize \emph{Grodno, 1866 r.}}};

\draw[-] (Walenty) -- (Helena)
    node[midway, above, draw=none, minimum width=0pt,
    font=\scriptsize] (Slub4_Gora) { }
    node[midway, below, draw=none, minimum width=0pt,
    font=\scriptsize] (Slub4) {\emph{Długa Kościelna, 1885 r.}};

\draw[->] (Slub3) -- (Slub4_Gora);

% --- POKOLENIE 5 ---
\node (StefanAleksander) [below=1.5cm of Slub4] {Stefan Aleksander Okrzeja \\
{\footnotesize \emph{Dębe, 1886 r.}} \\ BOHATER NARODOWY \flagaPL};

\draw[->] (Slub4) -- (StefanAleksander);

\end{tikzpicture}
\end{minipage}
\vspace*{0.75cm}
\end{figure}

Walenty Okrzeja\indeksosoby{Walenty}{Okrzeja}{1860}{Janówek} zmarł
w~1914~roku po~tym jak został wywieziony w~głąb Rosji. Jego żona
Helena\indeksosoby[Cieślińska]{Helena}{Okrzeja}{1866}{Grodno} dożyła
do~1949~roku, boleśnie przeżyła śmierć swoich czterech synów. Stan jej zdrowia
uległ pogorszeniu, gdy w grudniu 1944~roku niemiecki pocisk zniszczył jej dom
przy ulicy Grodzieńskiej w~Warszawie, w~którym mieszkała wraz z~córką
Kazimierą\footnote{Kazimiera Pelagia Mieszczanek z~domu Okrzeja opiekowała się
swoją matką Heleną do~ostatnich chwil jej życia.}.

\newpage
\BigBold{Potomkowie Sebastiana~i~Anny Piotrkowiczów}
\phantomsection\label{sec:sa_piotrkowicze}

Najmłodszy syn Bartłomieja Piotrkowicza,
Sebastian\indeksosoby{Sebastian}{Piotrkowicz}{1768}{Dębe Wielkie}, dnia
1~lutego~1789~roku w~Mińsku Mazowieckim wziął ślub z~Anną
Smugą\indeksosoby[Smuga]{Anna}{Piotrkowicz}{1770}{Dębe Wielkie} (ur.~1770~r.~--
zm.~1835~r.), dziewiętnastoletnią panną urodzoną w~rodzinnej wsi Piotrkowiczów.
Małżonkowie Piotrkowicze dochowali się wspólnie co~najmniej dziesięciorga
dzieci -- wszystkie przyszły na~świat w~Dębem.

\begin{figure}[!ht]
    \centering \includegraphics[width=0.75\linewidth]{
        1819_Michał_Piotrkowicz_Teresa_Kozłowska_akt_ślubu_parafia_Mińsk_Mazowiecki_wpis_33.jpg}
    \captionsetup{format=hang}
    \caption{Akt ślubu Michała Piotrkowicza i~Teresy Kozłowskiej -- par.
    Mińsk Mazowiecki 1819~rok (33/1819) \cite{par_minsk0}.}
    \label{fig:mpiotrkowicz_1819}
\end{figure}

Dziećmi Sebastiana\indeksosoby{Sebastian}{Piotrkowicz}{1768}{Dębe Wielkie}
i~Anny\indeksosoby[Smuga]{Anna}{Piotrkowicz}{1770}{Dębe Wielkie} Piotrkowiczów
były: Mateusz\indeksosoby{Mateusz}{Piotrkowicz}{1789}{Dębe Wielkie}
(ur.~1789~r.~-- zm.~1856~r.),
Agnieszka\indeksosoby{Agnieszka}{Piotrkowicz}{1790}{Dębe Wielkie}
(ur.~1790~r.~-- zm.~1814~r.),
Kazimierz\indeksosoby{Kazimierz}{Piotrkowicz}{1792}{Dębe Wielkie}
(ur.~1792~r.~-- zm.~1793~r.),
Maciej\indeksosoby{Maciej}{Piotrkowicz}{1793}{Dębe Wielkie}
(ur.~1793~r.~-- zm.~?),
Tomasz\indeksosoby{Tomasz}{Piotrkowicz}{1793}{Dębe Wielkie} (ur.~1793~r.~--
zm.~1796~r.), Marianna\indeksosoby{Marianna}{Piotrkowicz}{1797}{Dębe Wielkie}
(ur.~1797~r.~-- zm.~1810~r.),
Michał\indeksosoby{Michał}{Piotrkowicz}{1800}{Dębe Wielkie} (ur.~1800~r.~--
zm.~1856~r.), Józef\indeksosoby{Józef}{Piotrkowicz}{1804}{Dębe Wielkie}
(ur.~1804~r.~-- zm.~1864~r.),
Marianna\indeksosoby{Marianna}{Piotrkowicz}{1811}{Dębe Wielkie} (ur.~1811~r.~--
zm.~?) oraz Katarzyna\indeksosoby{Katarzyna}{Piotrkowicz}{1814}{Dębe Wielkie}
(ur.~1814~r.~-- zm.~?).

\begin{figure}[!ht]
    \vspace*{0.5cm}
    \centering \includegraphics[width=1.0\linewidth]{
        1825_Tomasz_Piotrkowicz_akt_chrztu_parafia_Mińsk_Mazowiecki_wpis_311.jpg}
    \captionsetup{format=hang}
    \caption{Akt chrztu Tomasza Piotrkowicza -- par. Mińsk Mazowiecki 1825~rok
    (311/1825) \cite{par_minsk1}.}
    \label{fig:tpiotrkowicz_1825}
\end{figure}

Z~perspektywy historii rodziny Gomulskich z~Desna najistotniejsi są~potomkowi
Michała Piotrkowicza\indeksosoby{Michał}{Piotrkowicz}{1800}{Dębe Wielkie},
gdyż był on~dziadkiem Jana
Piotrkowicza\indeksosoby{Jan}{Piotrkowicz}{1847}{Kąty Goździejewskie},
największego ze~współudziałowców spółki \emph{Desna}, o~którym szerzej pisałem
w~rozdziale V. Michał Piotrkowicz 14~lutego 1819~roku
w~Mińsku Mazowieckim wziął ślub z~Teresą
Kozłowską\indeksosoby[Kozłowska]{Teresa}{Piotrkowicz}{1799}{Dębe Wielkie}
(ur.~1799~r.~-- zm.~1858~r.) dwudziestoletnią panną z~Dębego -- patrz: ryc.
\ref{fig:mpiotrkowicz_1819}. Para dochowała się co~najmniej dziesięciorga
dzieci, wszystkie z~nich urodziły w~Dębem:
Józefy\indeksosoby{Józefa}{Piotrkowicz}{1820}{Dębe Wielkie} (ur.~1820~r.~--
zm.~?), Andrzeja\indeksosoby{Andrzej}{Piotrkowicz}{1823}{Dębe Wielkie}
(ur.~1823~r.~-- zm.~?),
Tomasza\indeksosoby{Tomasz}{Piotrkowicz}{1823}{Dębe Wielkie} (ur.~1825~r.~--
zm.~1900~r.), Józefa\indeksosoby{Józef}{Piotrkowicz}{1828}{Dębe Wielkie}
(ur.~1828~r.~-- zm.~1831~r.),
Karola\indeksosoby{Karol}{Piotrkowicz}{1830}{Dębe Wielkie} (ur.~1830~r.~--
zm.~1831~r.), Marianny\indeksosoby{Marianna}{Piotrkowicz}{1832}{Dębe Wielkie}
(ur.~1832~r.~-- zm.~?), Józefa\indeksosoby{Józef}{Piotrkowicz}{1834}{Dębe Wielkie}
(ur.~1834~r.~-- zm.~1864~r.),
Katarzyny\indeksosoby{Katarzyna}{Piotrkowicz}{1836}{Dębe Wielkie}
(ur.~1836~r.~-- zm.~?), Wojciecha\indeksosoby{Wojciech}{Piotrkowicz}{1839}{Dębe Wielkie}
(ur.~1839~r.~-- zm.~?) oraz Jakuba\indeksosoby{Jakub}{Piotrkowicz}{1840}{Dębe Wielkie}
(ur.~1840~r.~-- zm.~?).

\begin{figure}[!ht]
    \centering \includegraphics[width=0.78\linewidth]{
        1846_Tomasz_Piotrkowicz_Tekla_Wieczorek_akt_ślubu_parafia_Pustelnik_wpis_3.jpg}
    \captionsetup{format=hang}
    \caption{Akt ślubu Tomasza Piotrkowicza i~Tekli Wieczorek -- par.
    Pustelnik 1846~rok (3/1846) \cite{par_pustelnik}.}
    \label{fig:tpiotrkowicz_1846}
\end{figure}

Drugi syn Michała Piotrkowicza,
Tomasz\indeksosoby{Tomasz}{Piotrkowicz}{1823}{Dębe Wielkie}, był ojcem
Jana Piotrkowicza\indeksosoby{Jan}{Piotrkowicz}{1847}{Kąty Goździejewskie},
największego właściciela ziemskiego w~historii Desna oraz pierwszego męża
Karoliny
Gomulskiej\indeksosoby[Gomulska][Piotrkowicz]{Karolina}{Siporska}{1891}{Wólka Mińska},
córki Stanisława i~Julianny Gomulskich -- założycieli najmłodszej gałęzi
Gomulskich na~Deśnie. Tomasz
Piotrkowicz\indeksosoby{Tomasz}{Piotrkowicz}{1823}{Dębe Wielkie} 1~lutego
1846~roku ożenił się w~parafii pod wezwaniem św.~Katarzyny Aleksandryjskiej
w~Pustelniku z~Teklą
Wieczorek\indeksosoby[Wieczorek]{Tekla}{Piotrkowicz}{1826}{Kąty Goździejewskie}
(ur.~1826~r.~-- zm.~1897~r.) z~pobliskich Kątów Goździejewskich (patrz: ryc.
\ref{fig:tpiotrkowicz_1846}).

\begin{figure}[!ht]
    \vspace*{0.5cm}
    \centering \includegraphics[width=0.85\linewidth]{
        1847_Jan_Piotrkowicz_akt_chrztu_parafia_Pustelnik_wpis_4.jpg}
    \captionsetup{format=hang}
    \caption{Akt chrztu Jana Piotrkowicza -- par. Pustelnik 1847~rok (4/1847)
    \cite{par_pustelnik}.}
    \label{fig:jpiotrkowicz_1847}
\end{figure}

Tomasz Piotrkowicz\indeksosoby{Tomasz}{Piotrkowicz}{1823}{Dębe Wielkie}
i~Tekla Wieczorek\indeksosoby[Wieczorek]{Tekla}{Piotrkowicz}{1826}{Kąty Goździejewskie}
po~ślubie zamieszkali w~rodzinnej miejscowości panny młodej i~dochowali się
tam wspólnie co~najmniej dwanaściorga potomstwa:
Jana\indeksosoby{Jan}{Piotrkowicz}{1847}{Kąty Goździejewskie} (ur.~1847~r.~--
zm.~?), Julianny\indeksosoby{Julianna}{Piotrkowicz}{1849}{Kąty Goździejewskie}
(ur.~1849~r.~-- zm.~1919~r.),
Katarzyny\indeksosoby{Katarzyna}{Piotrkowicz}{1851}{Kąty Goździejewskie}
(ur.~1851~r.~-- zm.~?), Andrzeja\indeksosoby{Andrzej}{Piotrkowicz}{1853}{Kąty Goździejewskie}
(ur.~1853~r.~-- zm.~?), Łukasza\indeksosoby{Łukasz}{Piotrkowicz}{1855}{Kąty Goździejewskie}
(ur.~1855~r.~-- zm.~1927~r.),
Marianny\indeksosoby{Marianna}{Piotrkowicz}{1857}{Kąty Goździejewskie}
(ur.~1857~r.~-- zm.~1943~r.),
Józefa\indeksosoby{Józef}{Piotrkowicz}{1859}{Kąty Goździejewskie}
(ur.~1859~r.~-- zm.~1934~r.),
Małgorzaty\indeksosoby{Małgorzata}{Piotrkowicz}{1861}{Kąty Goździejewskie}
(ur.~1861~r.~-- zm.~1862~r.),
Franciszki\indeksosoby{Franciszka}{Piotrkowicz}{1864}{Kąty Goździejewskie}
(ur.~1864~r.~-- zm.~1928~r.),
Jakuba\indeksosoby{Jakub}{Piotrkowicz}{1866}{Kąty Goździejewskie}
(ur.~1866~r.~-- zm.~1867~r.),
Rozalii\indeksosoby{Rozalia}{Piotrkowicz}{1868}{Kąty Goździejewskie}
(ur.~1868~r.~-- zm.~1954~r.) oraz
Rocha\indeksosoby{Roch}{Piotrkowicz}{1870}{Kąty Goździejewskie}
(ur.~1870~r.~-- zm.~1884~r.).

\begin{figure}[!ht]
    \centering \includegraphics[width=0.78\linewidth]{
        1871_Jan_Piotrkowicz_Krystyna_Link_akt_ślubu_parafia_Długa_Kościelna_wpis_5.jpg}
    \captionsetup{format=hang}
    \caption{Akt ślubu Jana Piotrkowicza oraz Krystyny Link -- par. Długa
    Kościelna 1871~rok (5/1871) \cite{par_dluga}.}
    \label{fig:jpiotrkowicz_1871}
\end{figure}

\begin{figure}[!ht]
    \vspace*{0.5cm}
    \centering \includegraphics[width=0.98\linewidth]{
        1915_Karolina_Gomulska_Jan_Piotrkowicz_akt_ślubu_parafia_Długa_Kościelna_wpis_12.jpg}
    \captionsetup{format=hang}
    \caption{Akt ślubu Jana Piotrkowicza oraz Karoliny Gomulskiej -- par.
    Długa Kościelna 1915~rok (12/1915) \cite{par_dluga}.}
    \label{fig:kgomulska_1915_2}
\end{figure}

Najstarszy syn Tomasza i~Tekli Piotrkowiczów,
Jan\indeksosoby{Jan}{Piotrkowicz}{1847}{Kąty Goździejewskie}, 13~lutego
1871~roku wziął ślub z~Krystyną
Link\indeksosoby[Link]{Krystyna}{Piotrkowicz}{1845}{Olesin} (ur.~1845~r.~--
zm.~1915~r.) z~Kątów Mrowiskich (obecnie Olesin) -- patrz: ryc.
\ref{fig:jpiotrkowicz_1871}, gdzie zamieszkali po~ślubie i~wspólnie dochowali
się co~najmniej sześciorga dzieci:
Teofili\indeksosoby{Teofila}{Piotrkowicz}{1873}{Olesin} (ur.~1873~r.~-- zm.~?),
Franciszki\indeksosoby{Franciszek}{Piotrkowicz}{1875}{Olesin} (ur.~1875~r.~--
zm.~1959~r.),
Marcjanny\indeksosoby{Marcjanna}{Piotrkowicz}{1876}{Olesin} (ur.~1876~r.~--
zm.~1949~r.), Franciszki\indeksosoby{Franciszka}{Piotrkowicz}{1877}{Olesin}
(ur.~1879~r.~-- zm.~?), Andrzeja\indeksosoby{Andrzej}{Piotrkowicz}{1880}{Olesin}
(ur.~1880~r.~-- zm.~?) oraz
Stanisława\indeksosoby{Stanisław}{Piotrkowicz}{1883}{Olesin} (ur.~1883~r.~--
zm.~1957~r.). Kilka miesięcy po~śmierci swojej pierwszej żony, w~1915~roku,
68-letni wówczas Jan\indeksosoby{Jan}{Piotrkowicz}{1847}{Kąty Goździejewskie}
Piotrkowicz ożenił się z~dwudziestoczteroletnią Karoliną
Gomulską\indeksosoby[Gomulska][Piotrkowicz]{Karolina}{Siporska}{1891}{Wólka Mińska}
(ur.~1891~r.~-- zm.~1959~r.) -- patrz: ryc. \ref{fig:kgomulska_1915_2}. W~ten
sposób doszło do bezpośredniego połączenia rodzin Piotrkowiczów i~Gomulskich,
ponownego, bo~Jan i~Karolina byli ze~sobą spokrewnieni -- mieli wspólnego
przodka Bartłomieja
Piotrkowicza\indeksosoby{Bartłomiej}{Piotrkowicz}{1741}{Dębe Wielkie}, dla
Jana był to~prapradziadek, a~dla Karoliny praprapradziadek.

\begin{figure}[!ht]
\vspace*{1cm}
\centering
\hspace*{0.1cm}
\begin{tikzpicture}[
  % Definiujemy stałe odstępy
  node distance=1cm and 1cm, 
  every node/.style={draw, rectangle, align=center, minimum width=3.8cm,
  inner sep=5pt}]

  % --- POKOLENIE 0 ---
\node (Osiowa0) [draw=none, minimum width=0pt, inner sep=0pt] at (0,0) {};
\node (Bartlomiej) [left=1.5cm of Osiowa0] {Bartłomiej Piotrkowicz \\ {
    \footnotesize \emph{Dębe, 1741 r.}}};
\node (Zofia) [right=1.5cm of Osiowa0] {Zofia Kotczyk \\ {
    \footnotesize \emph{Dębe, 1743 r.}}};

\draw[-] (Bartlomiej) -- (Zofia)
    node[midway, below, draw=none, minimum width=0pt,
    font=\scriptsize] (Slub0) {\emph{Mińsk, 1759 r.}};

% STRZAŁKA W DÓŁ
\draw[->, thick] (Slub0.south) -- ++(0,-0.8);

\end{tikzpicture}
\end{figure}

\begin{figure}[!ht]
\centering
% --- PIERWSZY WYKRES ---
\begin{minipage}[t]{0.48\textwidth}
\centering
\begin{tikzpicture}[baseline=(Osiowa1.north), scale=0.7, transform shape,
    node distance=1cm and 1cm, 
    every node/.style={draw, rectangle, align=center, minimum width=3.5cm,
    inner sep=4pt}]

% --- POKOLENIE 1 ---
\node (Osiowa1) [draw=none, minimum width=0pt, inner sep=0pt] at (0,0) {};
\node (Sebastian) [left=1.5cm of Osiowa1] {Sebastian Piotrkowicz \\ {
    \footnotesize \emph{Dębe, 1768 r.}}};
\node (Anna) [right=1.5cm of Osiowa1] {Anna Smuga \\ {
    \footnotesize \emph{Dębe, 1770 r.}}};

\draw[-] (Sebastian) -- (Anna)
    node[midway, above, draw=none, minimum width=0pt,
    font=\scriptsize] (Slub1_Gora) { }
    node[midway, below, draw=none, minimum width=0pt,
    font=\scriptsize] (Slub1) {\emph{Mińsk, 1789 r.}};

% --- POKOLENIE 2 ---
\node (Osiowa2) [draw=none, minimum width=0pt, inner sep=0pt,
below=1.5cm of Slub1] {};
\node (Michal) [left=1.5cm of Osiowa2] {Michał Piotrkowicz \\ {
    \footnotesize \emph{Dębe, 1800 r.}}};
\node (Teresa) [right=1.5cm of Osiowa2] {Teresa Kozłowska \\ {
    \footnotesize \emph{Dębe, 1799 r.}}};
    
\draw[-] (Michal) -- (Teresa)
    node[midway, above, draw=none, minimum width=0pt,
    font=\scriptsize] (Slub2_Gora) { }
    node[midway, below, draw=none, minimum width=0pt,
    font=\scriptsize] (Slub2) {\emph{Mińsk, 1819 r.}};

\draw[->] (Slub1) -- (Slub2_Gora);

% --- POKOLENIE 3 ---
\node (Osiowa3) [draw=none, minimum width=0pt, inner sep=0pt,
below=1.5cm of  Slub2] {};
\node (Tomasz) [left=1.5cm of Osiowa3] {Tomasz Piotrkowicz \\ {
    \footnotesize \emph{Dębe, 1825 r.}}};
\node (Tekla) [right=1.5cm of Osiowa3] {Tekla Wieczorek \\ {
    \footnotesize \emph{Kąty Goździejewskie, 1826 r.}}};

\draw[-] (Tomasz) -- (Tekla)
    node[midway, above, draw=none, minimum width=0pt,
    font=\scriptsize] (Slub3_Gora) { }
    node[midway, below, draw=none, minimum width=0pt,
    font=\scriptsize] (Slub3) {\emph{Pustelnik, 1846 r.}};

\draw[->] (Slub2) -- (Slub3_Gora);

% --- POKOLENIE 4 ---
\node (Osiowa4) [draw=none, minimum width=0pt, inner sep=0pt,
below=1.5cm of  Slub3] {};
\node (Jan) [draw=none, left=1.5cm of Osiowa4] { \\
{\footnotesize \emph{ }}};
\node (Krystyna) [draw=none, right=1.5cm of Osiowa4] { \\ {
    \footnotesize \emph{ }}};

\draw[draw=none] (Jan) -- (Krystyna)
    node[midway, above, draw=none, minimum width=0pt,
    font=\scriptsize] (Slub4_Gora) { }
    node[midway, below, draw=none, minimum width=0pt,
    font=\scriptsize] (Slub4) {\emph{}};

\draw[->] (Slub3) -- (Slub4_Gora);
    
\end{tikzpicture}
\end{minipage}
\hfill % Doda elastyczny odstęp między wykresami
% --- DRUGI WYKRES ---
\begin{minipage}[t]{0.48\textwidth}
\centering
\begin{tikzpicture}[baseline=(Osiowa2.north), scale=0.7, transform shape,
    node distance=1cm and 1cm, 
    every node/.style={draw, rectangle, align=center, minimum width=3.5cm,
    inner sep=4pt}]

% --- POKOLENIE 1 ---
\node (Osiowa2) [draw=none, minimum width=0pt, inner sep=0pt] at (0,0) {};
\node (Pawel) [left=1.5cm of Osiowa2] {Paweł Piotrkowicz \\ {
    \footnotesize \emph{Dębe, 1760 r.}}};
\node (Dorota) [right=1.5cm of Osiowa2] {Dorota Okrzeja \\ {
    \footnotesize \emph{Dębe, 1781 r.}}};
    
\draw[-] (Pawel) -- (Dorota)
    node[midway, above, draw=none, minimum width=0pt,
    font=\scriptsize] (Slub2_Gora) { }
    node[midway, below, draw=none, minimum width=0pt,
    font=\scriptsize] (Slub2) {\emph{Mińsk, 1801 r.}};

% --- POKOLENIE 2 ---
\node (Osiowa3) [draw=none, minimum width=0pt, inner sep=0pt,
below=1.5cm of  Slub2] {};
\node (Tomasz) [left=1.5cm of Osiowa3] {Tomasz Piotrkowicz \\ {
    \footnotesize \emph{Dębe, 1804 r.}}};
\node (Marianna) [right=1.5cm of Osiowa3] {Marianna Lisiewska \\ {
    \footnotesize \emph{Dębe, 1811 r.}}};

\draw[-] (Tomasz) -- (Marianna)
    node[midway, above, draw=none, minimum width=0pt,
    font=\scriptsize] (Slub3_Gora) { }
    node[midway, below, draw=none, minimum width=0pt,
    font=\scriptsize] (Slub3) {\emph{Mińsk, 1829 r.}};

\draw[->] (Slub2) -- (Slub3_Gora);

% --- POKOLENIE 3 ---
\node (Osiowa4) [draw=none, minimum width=0pt, inner sep=0pt,
below=1.5cm of  Slub3] {};
\node (Agnieszka) [left=1.5cm of Osiowa4] {Agnieszka Piotrkowicz \\
{\footnotesize \emph{Budy Barcząckie, 1840 r.}}};
\node (Jan) [right=1.5cm of Osiowa4] {
    Jan Gomulski \\ {\footnotesize \emph{Zamienie, 1836 r.}}};

\draw[-] (Agnieszka) -- (Jan)
    node[midway, above, draw=none, minimum width=0pt,
    font=\scriptsize] (Slub4_Gora) { }
    node[midway, below, draw=none, minimum width=0pt,
    font=\scriptsize] (Slub4) {\emph{Mińsk, 1858 r.}};

\draw[->] (Slub3) -- (Slub4_Gora);

% --- POKOLENIE 4 ---
\node (Osiowa5) [draw=none, minimum width=0pt, inner sep=0pt,
below=1.5cm of Slub4] {};
\node (Stanislaw) [left=1.5cm of Osiowa5] {Stanisław Gomulski \\
{\footnotesize \emph{Wólka Mińska, 1861 r.}}};
\node (Julianna) [right=1.5cm of Osiowa5] {
    Julianna Gągol \\ {\footnotesize \emph{Dłużka, 1862 r.}}};

\draw[-] (Stanislaw) -- (Julianna)
    node[midway, above, draw=none, minimum width=0pt,
    font=\scriptsize] (Slub5_Gora) { }
    node[midway, below, draw=none, minimum width=0pt,
    font=\scriptsize] (Slub5) {\emph{Mińsk, 1881 r.}};

\draw[->] (Slub4) -- (Slub5_Gora);

% STRZAŁKA W DÓŁ
\draw[->] (Slub5.south) -- ++(0,-1.2);
    
\end{tikzpicture}
\end{minipage}
\end{figure}

\begin{figure}[!ht]
\centering
\hspace*{0.3cm}
\begin{tikzpicture}[
  % Definiujemy stałe odstępy
  node distance=1cm and 1cm, 
  every node/.style={draw, rectangle, align=center, minimum width=3.8cm,
  inner sep=5pt}]

% STRZAŁKA W DÓŁ
\draw[->, thick] (0, 0.8) -- (Osiowa0);

  % --- POKOLENIE 0 ---
\node (Osiowa0) [draw=none, minimum width=0pt, inner sep=0pt] at (0,0) {};
\node (Jan) [left=1.5cm of Osiowa0] {Jan Piotrkowicz \\ {
    \footnotesize \emph{Kąty Goździejewskie, 1847 r.}}};
\node (Karolina) [right=1.5cm of Osiowa0] {Karolina Gomulska \\ {
    \footnotesize \emph{Wólka Mińska, 1891 r.}}};

\draw[-] (Jan) -- (Karolina)
    node[midway, below, draw=none, minimum width=0pt,
    font=\scriptsize] (Slub0) {\emph{Długa Kościelna, 1915 r.}};

\end{tikzpicture}
\vspace*{0.5cm}
\end{figure}

Połączenie rodzin Piotrkowiczów i~Gomulskich miało też miejsce w~dalszych
pokoleniach, za~sprawą potomków najmłodszego syna Jana --
Stanisława\indeksosoby{Stanisław}{Piotrkowicz}{1883}{Olesin}. Dnia
24~lutego 1908~roku Stanisław Piotrkowicz wziął ślub w~Pustelniku z~Balbiną
Główką\indeksosoby[Główka]{Balbina}{Piotrkowicz}{1891}{Kąty Goździejewskie}
(ur.~1891~r.~-- zm.~1963~r.), panną z~Kątów Goździejewskich (patrz: ryc.
\ref{fig:jpiotrkowicz_1871}). Para po ślubie zamieszkała w~Olesinie, gdzie
dochowała się co~najmniej siedmiorga dzieci:
Piotra\indeksosoby{Piotr}{Piotrkowicz}{1911}{Olesin} (ur.~1911~r.~--
zm.~1911~r.), Heleny\indeksosoby{Helena}{Piotrkowicz}{1913}{Olesin}
(ur.~1913~r.~-- zm.~1980~r.), Józefa
Grzegorza\indeksosoby{Józef Grzegorz}{Piotrkowicz}{1915}{Olesin}
(ur.~1915~r.~-- zm.~1998~r.),
Stanisławy Magdaleny\indeksosoby[Piotrkowicz]{Stanisława Magdalena}{Kowalczyk}{1917}{Olesin}
(ur.~1917~r.~-- zm.~1987~r.), Jana\indeksosoby{Jan}{Piotrkowicz}{1919}{Olesin}
(ur.~1919~r.~-- zm.~1920~r.), Kazimierza\indeksosoby{Kazimierz}{Piotrkowicz}{1920}{Olesin}
(ur.~1920~r.~-- zm.~1999~r.) oraz Ryszarda\indeksosoby{Ryszard}{Piotrkowicz}{1927}{Olesin}
(ur.~1927~r.~-- zm.~2003~r.).

\begin{figure}[!ht]
    \vspace*{0.5cm}
    \centering \includegraphics[width=0.76\linewidth]{
        1883_Stanisław_Piotrkowicz_akt_chrztu_Długa_Kościelna_wpis_80.jpg}
    \captionsetup{format=hang}
    \caption{Akt chrztu Stanisława Piotrkowicza -- par. Długa Kościelna
    1883~rok (80/1883) \cite{par_dluga}.}
    \label{fig:spiotrkowicz_1883}
\end{figure}

\begin{figure}[!ht]
    \vspace*{1cm}
    \centering \includegraphics[width=0.76\linewidth]{
        1891_Balbina_Główka_Piotrkowicz_akt_chrztu_parafia_Pustelnik_wpis_24.jpg}
    \captionsetup{format=hang}
    \caption{Akt chrztu Balbiny Główki -- par. Pustelnik 1891~rok (24/1891)
    \cite{par_pustelnik}.}
    \label{fig:bglowka_1891}
\end{figure}

\begin{figure}[!ht]
    \vspace*{0.5cm}
    \centering \includegraphics[width=0.85\linewidth]{
        1908_Stanisław_Piotrkowicz_Balbina_Główka_akt_ślubu_parafia_Pustelnik_wpis_14.jpg}
    \captionsetup{format=hang}
    \caption{Akt ślubu Stanisława Piotrkowicza i~Balbiny Główki -- par.
    Pustelnik 1908~rok (14/1908) \cite{par_pustelnik}.}
    \label{fig:spiotrkowicz_1908}
\end{figure}

\begin{figure}[!ht]
    \centering \includegraphics[width=0.85\linewidth]{
        1911_Aleksander_Kowalczyk_akt_chrztu_parafia_Długa_Kościelna_wpis_120.jpg}
    \captionsetup{format=hang}
    \caption{Akt chrztu Aleksandra Kowalczyka -- par. Długa Kościelna 1911~rok
    (120/1911) \cite{par_dluga}.}
    \label{fig:akowalczyk_1911}
\end{figure}

Najmłodsza córka Stanisława\indeksosoby{Stanisław}{Piotrkowicz}{1883}{Olesin}
i~Balbiny\indeksosoby[Główka]{Balbina}{Piotrkowicz}{1891}{Kąty Goździejewskie}
Piotrkowiczów, Stanisława
Magdalena\indeksosoby[Piotrkowicz]{Stanisława Magdalena}{Kowalczyk}{1917}{Olesin}, dnia
4~lutego 1940~roku wzięła ślub w~parafii Świętych Apostołów Piotra
i~Pawła w~Dębem Wielkim z~Aleksandrem
Kowalczykiem\indeksosoby{Aleksander \flagaPL}{Kowalczyk}{1911}{Krzewina}\footnote{Co~ciekawe
Aleksander Kowalczyk działał w~Armii Krajowej pod pseudonimem
\enquote{Kosa} -- był członkiem w~stopniu kaprala 763~plutonu IV~kompanii
drugiego batalionu III~Rejonu \enquote{Dęby} (Rembertów) VII~Obwodu Armii
Krajowej \enquote{Obroża} (Warszawa-powiat) -- ryc. \ref{fig:kowalczyk_ak}.
Szczegółowy opis działalności 763~plutonu AK~można odnaleźć w~książce Jana
Kowalskiego, Ryszarda Rosia oraz Andrzeja Klimma pt. \emph{Czwarta Kompania
Rejon \enquote{Dęby} Armii Krajowej} \cite{kowalski}.}
dwudziestodziewięcioletnim kawalerem z~Krzewiny (ur.~1911~r.~-- zm.~1987~r.).
Para dochowała się wspólnie szóstki synów, jeden z~nich
Józef Franciszek\indeksosoby{Józef Franciszek}{Kowalczyk}{1953}{Krzewina}\footnote{Józef
Franciszek Kowalczyk jest ojcem chrzestnym autora książki.}
(ur.~1953~r.) ożenił się w~1977~roku w~Długiej Kościelnej z~córką
Stanisława\indeksosoby{Stanisław}{Gomulski}{1926}{Desno}
i~Jadwigi\indeksosoby[Poziemska]{Jadwiga Leokadia}{Gomulska}{1935}{Desno}
Gomulskich, Elżbietą
Marią\indeksosoby[Gomulska]{Elżbieta Maria}{Kowalczyk}{1954}{Desno} (ur.~1954~r.).
Józef i~Elżbieta Kowalczykowie są~ze~sobą bezpośrednio spokrewnieni -- mają
wspólnego przodka Bartłomieja
Piotrkowicza\indeksosoby{Bartłomiej}{Piotrkowicz}{1741}{Dębe Wielkie}, który
dla obojga z~nich jest 5x~pradziadkiem.

\begin{figure}[!ht]
    \vspace*{0.5cm}
    \centering \includegraphics[width=0.9\linewidth]{
        1917_Stanisława_Piotrkowicz_akt_chrztu_parafia_Dębe_Wielkie_wpis_17.jpg}
    \captionsetup{format=hang}
    \caption{Akt chrztu Stanisławy Magdaleny Piotrkowicz -- par. Dębe Wielkie
    1917~rok (17/1917) \cite{par_debe}.}
    \label{fig:smpiotrkowicz_1917}
    \vspace*{0.5cm}
\end{figure}

\begin{figure}[!ht]
    \vspace*{1.5cm}
    \centering \includegraphics[width=1.0\linewidth]{
        AK_Kowalczyk.jpg}
    \captionsetup{format=hang}
    \caption{Fragment książki \emph{Czwarta Kompania Rejon \enquote{Dęby}
    Armii Krajowej} przedstawiający członków AK działających na~terenie
    obecnej gminy Halinów \cite{kowalski}.}
    \label{fig:kowalczyk_ak}
\end{figure}

\newpage
\begin{figure}[!ht]
\centering
\hspace*{0.2cm}
\begin{tikzpicture}[
  % Definiujemy stałe odstępy
  node distance=1cm and 1cm, 
  every node/.style={draw, rectangle, align=center, minimum width=3.8cm,
  inner sep=5pt}]

  % --- POKOLENIE 0 ---
\node (Osiowa0) [draw=none, minimum width=0pt, inner sep=0pt] at (0,0) {};
\node (Bartlomiej) [left=1.5cm of Osiowa0] {Bartłomiej Piotrkowicz \\ {
    \footnotesize \emph{Dębe, 1741 r.}}};
\node (Zofia) [right=1.5cm of Osiowa0] {Zofia Kotczyk \\ {
    \footnotesize \emph{Dębe, 1743 r.}}};

\draw[-] (Bartlomiej) -- (Zofia)
    node[midway, below, draw=none, minimum width=0pt,
    font=\scriptsize] (Slub0) {\emph{Mińsk, 1759 r.}};

% STRZAŁKA W DÓŁ
\draw[->, thick] (Slub0.south) -- ++(0,-0.8);

\end{tikzpicture}
\end{figure}

\begin{figure}[!ht]
\centering
% --- PIERWSZY WYKRES ---
\begin{minipage}[t]{0.46\textwidth}
\centering
\begin{tikzpicture}[baseline=(Osiowa1.north), scale=0.7, transform shape,
    node distance=1cm and 1cm, 
    every node/.style={draw, rectangle, align=center, minimum width=3.5cm,
    inner sep=4pt}]

% --- POKOLENIE 1 ---
\node (Osiowa1) [draw=none, minimum width=0pt, inner sep=0pt] at (0,0) {};
\node (Sebastian) [left=1.5cm of Osiowa1] {Sebastian Piotrkowicz \\ {
    \footnotesize \emph{Dębe, 1768 r.}}};
\node (Anna) [right=1.5cm of Osiowa1] {Anna Smuga \\ {
    \footnotesize \emph{Dębe, 1770 r.}}};

\draw[-] (Sebastian) -- (Anna)
    node[midway, above, draw=none, minimum width=0pt,
    font=\scriptsize] (Slub1_Gora) { }
    node[midway, below, draw=none, minimum width=0pt,
    font=\scriptsize] (Slub1) {\emph{Mińsk, 1789 r.}};

% --- POKOLENIE 2 ---
\node (Osiowa2) [draw=none, minimum width=0pt, inner sep=0pt,
below=1.5cm of Slub1] {};
\node (Michal) [left=1.5cm of Osiowa2] {Michał Piotrkowicz \\ {
    \footnotesize \emph{Dębe, 1800 r.}}};
\node (Teresa) [right=1.5cm of Osiowa2] {Teresa Kozłowska \\ {
    \footnotesize \emph{Dębe, 1799 r.}}};
    
\draw[-] (Michal) -- (Teresa)
    node[midway, above, draw=none, minimum width=0pt,
    font=\scriptsize] (Slub2_Gora) { }
    node[midway, below, draw=none, minimum width=0pt,
    font=\scriptsize] (Slub2) {\emph{Mińsk, 1819 r.}};

\draw[->] (Slub1) -- (Slub2_Gora);

% --- POKOLENIE 3 ---
\node (Osiowa3) [draw=none, minimum width=0pt, inner sep=0pt,
below=1.5cm of  Slub2] {};
\node (Tomasz) [left=1.5cm of Osiowa3] {Tomasz Piotrkowicz \\ {
    \footnotesize \emph{Dębe, 1825 r.}}};
\node (Tekla) [right=1.5cm of Osiowa3] {Tekla Wieczorek \\ {
    \footnotesize \emph{Kąty Goździejewskie, 1826 r.}}};

\draw[-] (Tomasz) -- (Tekla)
    node[midway, above, draw=none, minimum width=0pt,
    font=\scriptsize] (Slub3_Gora) { }
    node[midway, below, draw=none, minimum width=0pt,
    font=\scriptsize] (Slub3) {\emph{Pustelnik, 1846 r.}};

\draw[->] (Slub2) -- (Slub3_Gora);

% --- POKOLENIE 4 ---
\node (Osiowa4) [draw=none, minimum width=0pt, inner sep=0pt,
below=1.5cm of  Slub3] {};
\node (Jan) [left=1.5cm of Osiowa4] {Jan Piotrkowicz \\ {
    \footnotesize \emph{Kąty Goździejewskie, 1847 r.}}};
\node (Krystyna) [right=1.5cm of Osiowa4] {Krystyna Link \\ {
    \footnotesize \emph{Kąty Mrowiskie, 1845 r.}}};

\draw[-] (Jan) -- (Krystyna)
    node[midway, above, draw=none, minimum width=0pt,
    font=\scriptsize] (Slub4_Gora) { }
    node[midway, below, draw=none, minimum width=0pt,
    font=\scriptsize] (Slub4) {\emph{Pustelnik, 1871 r.}};

\draw[->] (Slub3) -- (Slub4_Gora);

% --- POKOLENIE 5 ---
\node (Osiowa5) [draw=none, minimum width=0pt, inner sep=0pt,
below=1.5cm of  Slub4] {};
\node (Stanislaw) [left=1.5cm of Osiowa5] {Stanisław Piotrkowicz \\ {
    \footnotesize \emph{Olesin, 1883 r.}}};
\node (Balbina) [right=1.5cm of Osiowa5] {Balbina Główka \\ {
    \footnotesize \emph{Kąty Goździejewskie, 1891 r.}}};

\draw[-] (Stanislaw) -- (Balbina)
    node[midway, above, draw=none, minimum width=0pt,
    font=\scriptsize] (Slub5_Gora) { }
    node[midway, below, draw=none, minimum width=0pt,
    font=\scriptsize] (Slub5) {\emph{Pustelnik, 1908 r.}};

\draw[->] (Slub4) -- (Slub5_Gora);

% --- POKOLENIE 6 ---
\node (Osiowa6) [draw=none, minimum width=0pt, inner sep=0pt,
below=1.5cm of  Slub5] {};
\node (Stanislawa) [left=1.5cm of Osiowa6] {
    Stanisława Magdalena \\ Piotrkowicz \\ {\footnotesize \emph{
        Olesin, 1917 r.}}};
\node (Aleksander) [right=1.5cm of Osiowa6] {Aleksander Kowalczyk \\ {
    \footnotesize \emph{Krzewina, 1911 r.}} \\ ŻOŁNIERZ AK \flagaPL};

\draw[-] (Stanislawa) -- (Aleksander)
    node[midway, above, draw=none, minimum width=0pt,
    font=\scriptsize] (Slub6_Gora) { }
    node[midway, below, draw=none, minimum width=0pt,
    font=\scriptsize] (Slub6) {\emph{Dębe Wielkie, 1940 r.}};

\draw[->] (Slub5) -- (Slub6_Gora);

% STRZAŁKA W DÓŁ
\draw[->] (Slub6.south) -- ++(0,-1.2);
    
\end{tikzpicture}
\end{minipage}
\hfill % Doda elastyczny odstęp między wykresami
% --- DRUGI WYKRES ---
\begin{minipage}[t]{0.46\textwidth}
\centering
\begin{tikzpicture}[baseline=(Osiowa2.north), scale=0.7, transform shape,
    node distance=1cm and 1cm, 
    every node/.style={draw, rectangle, align=center, minimum width=3.5cm,
    inner sep=4pt}]

% --- POKOLENIE 1 ---
\node (Osiowa2) [draw=none, minimum width=0pt, inner sep=0pt] at (0,0) {};
\node (Pawel) [left=1.5cm of Osiowa2] {Paweł Piotrkowicz \\ {
    \footnotesize \emph{Dębe, 1760 r.}}};
\node (Dorota) [right=1.5cm of Osiowa2] {Dorota Okrzeja \\ {
    \footnotesize \emph{Dębe, 1781 r.}}};
    
\draw[-] (Pawel) -- (Dorota)
    node[midway, above, draw=none, minimum width=0pt,
    font=\scriptsize] (Slub2_Gora) { }
    node[midway, below, draw=none, minimum width=0pt,
    font=\scriptsize] (Slub2) {\emph{Mińsk, 1801 r.}};

% --- POKOLENIE 2 ---
\node (Osiowa3) [draw=none, minimum width=0pt, inner sep=0pt,
below=1.5cm of  Slub2] {};
\node (Tomasz) [left=1.5cm of Osiowa3] {Tomasz Piotrkowicz \\ {
    \footnotesize \emph{Dębe, 1804 r.}}};
\node (Marianna) [right=1.5cm of Osiowa3] {Marianna Lisiewska \\ {
    \footnotesize \emph{Dębe, 1811 r.}}};

\draw[-] (Tomasz) -- (Marianna)
    node[midway, above, draw=none, minimum width=0pt,
    font=\scriptsize] (Slub3_Gora) { }
    node[midway, below, draw=none, minimum width=0pt,
    font=\scriptsize] (Slub3) {\emph{Mińsk, 1829 r.}};

\draw[->] (Slub2) -- (Slub3_Gora);

% --- POKOLENIE 3 ---
\node (Osiowa4) [draw=none, minimum width=0pt, inner sep=0pt,
below=1.5cm of  Slub3] {};
\node (Agnieszka) [left=1.5cm of Osiowa4] {Agnieszka Piotrkowicz \\
{\footnotesize \emph{Budy Barcząckie, 1840 r.}}};
\node (Jan) [right=1.5cm of Osiowa4] {
    Jan Gomulski \\ {\footnotesize \emph{Zamienie, 1836 r.}}};

\draw[-] (Agnieszka) -- (Jan)
    node[midway, above, draw=none, minimum width=0pt,
    font=\scriptsize] (Slub4_Gora) { }
    node[midway, below, draw=none, minimum width=0pt,
    font=\scriptsize] (Slub4) {\emph{Mińsk, 1858 r.}};

\draw[->] (Slub3) -- (Slub4_Gora);

% --- POKOLENIE 4 ---
\node (Osiowa5) [draw=none, minimum width=0pt, inner sep=0pt,
below=1.5cm of Slub4] {};
\node (Stanislaw) [left=1.5cm of Osiowa5] {Stanisław Gomulski \\
{\footnotesize \emph{Wólka Mińska, 1861 r.}}};
\node (Julianna) [right=1.5cm of Osiowa5] {
    Julianna Gągol \\ {\footnotesize \emph{Dłużka, 1862 r.}}};

\draw[-] (Stanislaw) -- (Julianna)
    node[midway, above, draw=none, minimum width=0pt,
    font=\scriptsize] (Slub5_Gora) { }
    node[midway, below, draw=none, minimum width=0pt,
    font=\scriptsize] (Slub5) {\emph{Mińsk, 1881 r.}};

\draw[->] (Slub4) -- (Slub5_Gora);

% --- POKOLENIE 5 ---
\node (Osiowa6) [draw=none, minimum width=0pt, inner sep=0pt,
below=1.5cm of Slub5] {};
\node (Jozef) [left=1.5cm of Osiowa6] {Józef Gomulski \\
{\footnotesize \emph{Wólka Mińska, 1895 r.}}};
\node (Feliksa) [right=1.5cm of Osiowa6] {
    Feliksa Tkaczyk \\ {\footnotesize \emph{Desno, 1897 r.}}};

\draw[-] (Jozef) -- (Feliksa)
    node[midway, above, draw=none, minimum width=0pt,
    font=\scriptsize] (Slub6_Gora) { }
    node[midway, below, draw=none, minimum width=0pt,
    font=\scriptsize] (Slub6) {\emph{Długa Kościelna, 1916 r.}};

\draw[->] (Slub5) -- (Slub6_Gora);

% --- POKOLENIE 6 ---
\node (Osiowa7) [draw=none, minimum width=0pt, inner sep=0pt,
below=1.5cm of Slub6] {};
\node (Stanislaw2) [left=2cm of Osiowa7] {Stanisław Gomulski \\
{\footnotesize \emph{Desno, 1926 r.}}};
\node (Jadwiga) [right=2cm of Osiowa7] {
    Jadwiga Leokadia \\ Poziemska \\ {\footnotesize \emph{Desno, 1935 r.}}};

\draw[-] (Stanislaw2) -- (Jadwiga)
    node[midway, above, draw=none, minimum width=0pt,
    font=\scriptsize] (Slub7_Gora) { }
    node[midway, below, draw=none, minimum width=0pt,
    font=\scriptsize] (Slub7) {\emph{Długa Kościelna, 1951 r.}};

\draw[->] (Slub6) -- (Slub7_Gora);

% STRZAŁKA W DÓŁ
\draw[->] (Slub7.south) -- ++(0,-1.2);
    
\end{tikzpicture}
\end{minipage}
\end{figure}

\begin{figure}[!ht]
\centering
\hspace*{0.7cm}
\begin{tikzpicture}[
  % Definiujemy stałe odstępy
  node distance=1cm and 1cm,
  every node/.style={draw, rectangle, align=center, minimum width=3.8cm,
  inner sep=5pt}]

% STRZAŁKA W DÓŁ
\draw[->, thick] (0, 0.8) -- (Osiowa0);

  % --- POKOLENIE 0 ---
\node (Osiowa0) [draw=none, minimum width=0pt, inner sep=0pt] at (0,0) {};
\node (Jozef) [left=1.5cm of Osiowa0] {Józef Franciszek \\ Kowalczyk \\ {
    \footnotesize \emph{Krzewina, 1953 r.}}};
\node (Elzbieta) [right=1.5cm of Osiowa0] {Elżbieta Maria \\ Gomulska \\ {
    \footnotesize \emph{Desno, 1954 r.}}};

\draw[-] (Jozef) -- (Elzbieta)
    node[midway, below, draw=none, minimum width=0pt,
    font=\scriptsize] (Slub0) {\emph{Długa Kościelna, 1977 r.}};

\end{tikzpicture}
\vspace*{0.5cm}
\end{figure}

Potomkowie Piotrkowiczów i~Okrzejów żyją do~dnia dzisiejszego na~terenie
obecnych gmin Dębe Wielkie oraz Halinów, część z~nich nosi nadal nazwisko
swoich przodków mieszkając w~Olesinie, Cisiu czy Dębem Wielkim. Poza rodziną
Gomulskich z~Desna, spokrewnione są z~nimi również inne liczne rodziny
zamieszkujące okoliczne miejscowości~-- m.in.~Książkowie, Mirosowie oraz
Malesowie z~Olesina, Wichorowscy oraz Bokowie z~Desna, Kowalczykowie
z~Krzewiny czy Grubowie z~Dębego Wielkiego.